% 型付き AI エージェント言語
% AIPL: Actor based Intelligence Parallel Language — 設計から実現 —（第5版）
%   lualatex -interaction=nonstopmode aipl_paper5_ja.tex   （3回）
\documentclass[11pt,a4paper]{ltjsarticle}
\usepackage{amsmath,amssymb}
\usepackage{amsthm}
\usepackage{listings}
\usepackage{xcolor}
\usepackage{geometry}
\usepackage{array}
\usepackage{booktabs}
\usepackage{longtable}
\usepackage[hidelinks]{hyperref}
\geometry{margin=21mm}

\newtheorem{theorem}{定理}
\newtheorem{lemma}[theorem]{補題}
\newtheorem{corollary}[theorem]{系}
\theoremstyle{definition}
\newtheorem{definition}[theorem]{定義}
\theoremstyle{remark}
\newtheorem{remark}[theorem]{注意}
\newtheorem{finding}[theorem]{所見}

\definecolor{codebg}{gray}{0.96}
\lstset{
  basicstyle=\ttfamily\small, backgroundcolor=\color{codebg},
  frame=single, framesep=4pt, rulecolor=\color{gray},
  breaklines=true, showstringspaces=false,
  columns=fullflexible, keepspaces=true,
  keywordstyle=\bfseries, commentstyle=\itshape\color{gray!130},
  xleftmargin=0pt, xrightmargin=0pt,
}
\lstdefinelanguage{AIPL}{
  morekeywords={class,method,var,new,now,future,await,reply,send,
    timeout,else,if,while,do,select,case,self,sender,remote,print,call,
    int,float,string,bool,unit,any,true,false},
  sensitive=true, morecomment=[l]{//}, morestring=[b]",
}
\lstnewenvironment{aipl}[1][]{\lstset{language=AIPL,#1}}{}
\lstnewenvironment{term}[1][]{\lstset{language=,basicstyle=\ttfamily\footnotesize,#1}}{}

\title{\bfseries 型付き AI エージェント言語\\[1.5mm]
       \bfseries AIPL: Actor based Intelligence Parallel Language\\[2mm]
       \large --- 設計から実現 ---\\[1mm]
       \normalsize 第 5 版：現在の仕様と、二つの源流からの採否}
\author{児玉靖司\\\small 法政大学経営学部\\\small \texttt{yass@hosei.ac.jp}}
\date{2026 年 8 月 23 日（第 5 版）}

\begin{document}
\maketitle

\begin{abstract}
AI エージェントは、外部のモデルを呼び、他のエージェントを待ち、道具を掴んで放す ---
つまり\textbf{並行で、副作用を持ち、いつ返るか分からない}プログラムである。
にもかかわらず、いま広く使われている書き方には型が無い。
何を待っているのか、何を汚すのか、返らなかったらどうなるのかが、
プログラムの表面に現れない。

本稿は、この隙間を埋めるために作った言語 AIPL（Actor based Intelligence
Parallel Language、処理系名 ABCLc+）の\textbf{現在の仕様}と、
その設計判断をまとめる。
AIPL はアクターを単位とする\textbf{型付き並行言語}であり、
一つのエージェントは一つのアクター、役割どうしのやり取りはメッセージ送信である。

設計は二つの源流から出発している。
\textbf{ABCL/1}（Yonezawa, Briot, Shibayama, OOPSLA 1986）と、
\textbf{小林直樹}らの並行計算のための型システムの系譜である。
本稿は、両者の機構をひとつずつ挙げ、
\textbf{採ったか落としたか・その理由・メリット・デメリット}を、
すべて動く最小プログラムとともに示す。
ABCL/1 からは三種のメッセージ送信・オブジェクト内部の逐次性・\texttt{reply} による返信・
返信先の一級性を採り、\textit{express mode} による割り込みと
\texttt{script} の \texttt{where} ガードを落とした。
小林の系譜からは資源使用解析（効果）・線形型・デッドロック自由な型システム・
サイズ型（粗い代用として期限）を採り、
交差型による高階モデル検査・多相メソッド・完全な所有権型を落とした。

結果として、現在の型システムが押さえるものは七つである。
\textbf{(1)} \textbf{効果} --- \texttt{ai net io mut time mem fs log} の 8 種を
本体から集め、\texttt{now}/\texttt{await} で呼び先から伝播させる。
モデルを呼ぶメソッドは \texttt{ai} を隠せない。
\textbf{(2)} \textbf{期限} --- 待ちには \texttt{timeout} を要求する。
\textbf{(3)} \textbf{失敗の型} --- \texttt{else} を書かない待ちは
\texttt{result}$\langle\tau\rangle$ を返し、確かめずには使えない。
\textbf{(4)} \textbf{返信先の一級性} --- 線形型で「ちょうど一度」を保ったまま委譲できる。
\textbf{(5)} \textbf{資源の順序} --- \texttt{acquire}/\texttt{release} の対に加え、
取得の入れ子から\textbf{全体順序}を読み取り、逆順の取得を止める。
\textbf{(6)} \textbf{義務レベル} --- 待ちの循環を止める。注釈でも書けるし、書かなければ推論される。
\textbf{(7)} \textbf{セッション型} --- 役割どうしのやり取りの順序を静的に照合する。
同期の呼び先は\textbf{展開}するので、セッションがアクターをまたいでも追える。
併せて\textbf{メッシュ配備}を持ち、アクターのソースを他ノードへ送って
相手先で構文解析・型検査してから実体化する。効果注釈がコードと一緒に運ばれる。

\textbf{デッドロック}については、AIPL で待ちが返らなくなる経路が\textbf{四つ}あり、
一つのものとして扱うと答えを誤ることを示す ---
循環待ち、返信されない待ち、\texttt{select} が来ないメッセージを待つ場合、
そして資源を逆順に取る場合である。
四つとも型で押さえた結果、\textbf{待ちの宛先がクラスとして分かるか、
宛先ノードのレベルが宣言されている範囲では、デッドロック自由が型で保証される}。
成立の条件と、依然として残る範囲を明記する。

三役（立案・求解・査読）の AI エージェントを題材に、
型が何を止めるかを実測で示す --- 期限を外すと警告が出て後続の型が合わなくなり、
\texttt{ai} 効果を隠すと弾かれ、役割の順序を入れ替えるとセッション型が弾き、
失敗を確かめずに次へ渡すと型が合わない。
最新仕様の代表的なサンプル 20 本を、コード・出力・メリット・デメリットとともに解説し、
そのうち 18 本が\textbf{三つの独立な処理系で出力一致}することを確認した。
弾かれるべきプログラムは 31 本ある。
\end{abstract}

\tableofcontents
\newpage

% =====================================================================
\section{はじめに}
\label{sec:intro}

\subsection{AI エージェントには型が無い}

いま「AI エージェント」と呼ばれているプログラムの多くは、
次のような形をしている --- ある役割がモデルを呼び、
その答えを別の役割へ渡し、道具を掴んで結果を書き戻す。
役割は並行に動き、互いの返事を待つ。

これは新しい種類のプログラムではない。
\textbf{アクターとメッセージ}そのものである。
違うのは、待ちの相手が\textbf{機外のモデル}であり、
いつ返るか、そもそも返るかが分からないことだけである。

にもかかわらず、この種のプログラムには型が無い。
役割の呼び出しは文字列の受け渡しで、
そのメソッドが\textbf{モデルを呼ぶのか}、
\textbf{どれだけ待つのか}、
\textbf{返らなかったら何が起きるのか}が、どこにも書かれていない。
順序の約束 --- 立案してから解き、解いてから査読する --- も、
書いた人の頭の中にしかない。

失敗はモデルの間違いより先に、この隙間で起きる。
査読役が求解役を待ち、求解役が査読役を待って止まる。
期限を書き忘れた待ちが返らない。
時間切れの空文字列が、正常な答えのふりをして次の役割へ流れていく。
\textbf{どれも型の話である}。

\subsection{二つの世界を混ぜない}

AIPL は二つの目的のために作られている。
\textbf{分散 AI エージェント}の記述と、\textbf{実時間 OS} の記述である。
前者は「機外のモデルを呼ぶので、いつ返るか分からない」世界であり、
後者は「期限を守れないことが故障である」世界である。
両方をアクターとメッセージで書きたい。
そして\textbf{両者を混ぜてはいけない}。
実時間で動くアクターが機外のモデルを同期で待ってはならない。
これを規約ではなく型検査で禁じたい --- これが設計を貫く動機である。

この一文が、本稿に出てくる仕組みのほとんどを説明する。
「モデルを呼ぶ」を型に出すのが\textbf{効果}であり、
「いつ返るか分からない」を型に出すのが\textbf{期限}と
\textbf{\texttt{result}$\langle\tau\rangle$}であり、
「実時間側が上位を待たない」を型に出すのが\textbf{義務レベル}である。

\subsection{二つの源流}

言語を一から作ってはいない。二つの先行研究から出発している。

\paragraph{ABCL/1（1986）。}
Yonezawa, Briot, Shibayama によるアクター指向の並行オブジェクト言語である。
中核は\textbf{三種のメッセージ送信}（past / now / future）で、
これはそのまま「待たない／待つ／後で受け取る」という区別に対応する。
オブジェクトは一度に一つのメッセージしか処理せず、
\textit{express mode} という割り込み経路と、
\texttt{script} の \texttt{where} ガードによる受理条件を持つ。
ただし\textbf{動的型}であり、型システムを持たない。

\paragraph{小林らの型システム（1997---）。}
ABCL/1 が持たなかった基礎づけを与える系譜である。
$\pi$ 計算に対する線形型、資源使用解析、デッドロック自由な型システム、
サイズ型・精製型、交差型による高階モデル検査などが含まれる。
「通信路をちょうど一度使う」「資源を取ったら放す」「待ちに順序を課す」
といった性質を、実行前に型で言うための道具立てである。

\paragraph{本稿の立場。}
AIPL は、\textbf{ABCL/1 の書き味に小林の型システムを載せる}試みである。
ただし両者をそのまま合わせたのではない。
実際に処理系を三つ書き、582 本のプログラムを通しながら、
機構ごとに採否を決めてきた。
本稿はその\textbf{結果}を、機構ごとの理由・メリット・デメリットとして述べる。

\subsection{本稿の構成}

\S\ref{sec:spec} に AIPL の現在の姿を一望できる形で置く。
\S\ref{sec:abcl} が ABCL/1 の機能の採否、
\S\ref{sec:koba} が小林らの型システムの機能の採否、
\S\ref{sec:orig} がどちらの源流にも無い三つの機構である。
\S\ref{sec:table} に採否の対応表を置く。
\S\ref{sec:dl} がデッドロックの四つの経路とそれぞれの道具、
\S\ref{sec:agent} が\textbf{型付き AI エージェント}を題材にした実測、
\S\ref{sec:samples} が代表的なサンプル 20 本、
\S\ref{sec:impl} が三つの処理系と検証、
\S\ref{sec:discuss} が考察である。

本稿のプログラムはすべて実際に処理系へ通した。
出力は貼り付けたものではなく、実行して得たものである。

% =====================================================================
\section{AIPL の現在の姿}
\label{sec:spec}

採否の議論に入る前に、いまの言語を一望しておく。
以下はすべて現在の仕様であり、実装も三つある（\S\ref{sec:impl}）。

\subsection{アクターとメッセージ}

プログラムは\textbf{クラス}と\textbf{トップレベルの文}からなる。
クラスの実体が\textbf{アクター}で、状態（フィールド）とメソッドを持つ。
アクターは一度に一つのメッセージしか処理しない ---
だから状態を触るのに排他制御を書かない。

送信は三種ある。

\begin{center}
\begin{tabular}{lll}
\toprule
書き方 & 意味 & 返り値 \\
\midrule
\texttt{send a.m(x);} & 送って待たない & 無し \\
\texttt{now a.m(x) timeout $n$ else $e$} & 返信まで待つ & $\tau$ \\
\texttt{now a.m(x) timeout $n$} & 同上（\texttt{else} 無し） & \texttt{result}$\langle\tau\rangle$ \\
\texttt{future a.m(x)} $\to$ \texttt{await $f$} & 後で受け取る & \texttt{future}$\langle\tau\rangle$ \\
\bottomrule
\end{tabular}
\end{center}

返信は \texttt{reply(e)} で行い、\textbf{高々一度}である。
返信先そのものを値として取り出すこともできる（\texttt{replyto}）。

\subsection{型・効果・期限}

型は Hindley--Milner 流に推論する。
メソッドの戻り値型は、書かなければ \texttt{reply} から推論される
（内部的にはメソッドごとに戻り値型変数 $\rho$ を置く）。
書けば、その型が正となり、全経路での \texttt{reply} が課される。

判断は三つ組である。

\[
\Gamma \vdash e : \tau \,!\, \varepsilon
\]

$\varepsilon$ が\textbf{効果}で、次の 8 種からなる。

\begin{center}
\begin{tabular}{ll}
\toprule
効果 & 意味 \\
\midrule
\texttt{ai} & 機外のモデルを呼ぶ \\
\texttt{net} & ネットワークに触る \\
\texttt{io} & 装置に出力する \\
\texttt{mut} & 自分の状態を変える \\
\texttt{time} & 時間を消費する（\texttt{wait}） \\
\texttt{mem} & 永続記憶に触る \\
\texttt{fs} & ファイルに触る \\
\texttt{log} & 記録を残す（\texttt{print}） \\
\bottomrule
\end{tabular}
\end{center}

効果は本体から集められ、\texttt{now} / \texttt{await} で\textbf{呼び先から伝播する}。
注釈 \texttt{!\{...\}} を書けば、本体の効果がそこに収まることが検査される。
書かなければ推論に任せる（gradual）。

待ちには\textbf{期限}を課す。
\texttt{timeout $n$ else $e$} は $n$ ミリ秒待って諦め $e$ を返す。
\texttt{else} を書かなければ \texttt{result}$\langle\tau\rangle$ が返り、
\texttt{is\_ok} / \texttt{value} で確かめないと使えない。

\subsection{七つの規律}

型と効果と期限のうえに、次の七つが載っている。

\begin{center}
\small
\begin{tabular}{p{0.20\textwidth}p{0.34\textwidth}p{0.36\textwidth}}
\toprule
機構 & 書き方 & 止めるもの \\
\midrule
効果 & \texttt{!\{ai, net\}} & 実時間側が機外のモデルを待つこと \\
期限 & \texttt{timeout $n$ else $e$} & 返らない待ち \\
失敗の型 & \texttt{timeout $n$}（\texttt{else} 無し） & 期限切れの値を正常値として流すこと \\
返信先の一級性 & \texttt{replyto} / \texttt{answer} / 引数型 \texttt{reply} & 返信の取りこぼしと二重返信 \\
資源の順序 & \texttt{acquire} / \texttt{release}、\texttt{resource\_order} & 放し忘れ・二重取得・逆順の取得 \\
義務レベル & \texttt{@$n$}（書かなければ推論） & 待ちの循環 \\
セッション型 & \texttt{protocol\_define} / \texttt{\_start} / \texttt{\_end} & やり取りの順序違反 \\
\bottomrule
\end{tabular}
\end{center}

加えて\textbf{メッシュ配備}がある ---
アクターのソースを他ノードへ送り、相手先で構文解析・型検査してから実体化する。
効果注釈がコードと一緒に運ばれ、受け入れ側の方針と照合される。

\subsection{逃げ道}

既定はすべて厳しい側に倒してある。
緩める向きの環境変数を用意しているが、既定で使うことは想定していない。

\begin{term}
AIOS_LAX_WAIT=1        ... 循環待ちを警告に降格
AIOS_LAX_ACTOR=1       ... アクター型を区別しない
AIOS_LAX_ASSIGN=1      ... 未宣言の名前への代入を許す
AIOS_STRICT_DEADLINE=1 ... 期限のない待ちをエラーに昇格
AIOS_STRICT_PROTOCOL=1 ... 追えないセッションのやり残しを警告
AIOS_STRICT_RESOURCE=1 ... 名前がリテラルでない acquire の場所を知らせる
AIOS_SHOW_LEVELS=1     ... 推論した義務レベルと資源の全体順序を見せる
\end{term}

% =====================================================================
\section{ABCL/1 の機能と、その採否}
\label{sec:abcl}

ABCL/1 が持っていた機構をひとつずつ挙げ、
AIPL がそれを採ったか落としたかを、理由・メリット・デメリットとともに述べる。

\subsection{三種のメッセージ送信 ---【採る】}


最も本質的な継承である。

\begin{center}
\begin{tabular}{llll}
\toprule
AIPL の書き方 & ABCL/1 & 意味 & 式か文か \\
\midrule
\texttt{send t.m(a);}                          & past   & 送って待たない & 文 \\
\texttt{now t.m(a) timeout $n$ else $e$}       & now    & 送って返信を待つ & 式 \\
\texttt{future t.m(a)} / \texttt{await f ...}  & future & 引換券を受け取る & 式 \\
\bottomrule
\end{tabular}
\end{center}

\begin{aipl}
class Calc {
  method twice(x) : int { reply(x * 2); }
  method slow(x)  : int { reply(x + 100); }
}
var c = new Calc();
send c.twice(1);                                            // 待たない
print("now    = " ++ (now c.twice(21) timeout 100 else 0)); // 待つ
var f = future c.slow(5);                                   // 引換券
print("future = " ++ (await f timeout 100 else 0));         // あとで受け取る
\end{aipl}
\begin{term}
now    = 42
future = 105
\end{term}

\paragraph{なぜ採ったか。}
この三分割は、そのまま AIPL が必要とする区別である。
実時間側は「待たない」を多用し、エージェント側は「後で受け取る」を使う。
ABCL/f は future だけを残したが、AIPL は三種すべてを残した。
そのぶん型付けは面倒になる --- \texttt{send} は返信を受け取らないので戻り値型を要求せず、
\texttt{now} は式なのでその場に戻り値型が要り、
\texttt{future} は $\tau$ を返すメソッドに \texttt{future}$\langle\tau\rangle$ を与える。
この面倒は、三種の区別が設計上必要だという判断に対する代価である。

\paragraph{採った理由。}
「待たない／待つ／後で受け取る」は、並行プログラムを書くうえで
\textbf{どうしても要る区別}である。
そしてこの三つは型のうえでも別物になる ---
\texttt{send} は値を返さず、\texttt{now} は $\tau$ を返し、
\texttt{future} は \texttt{future}$\langle\tau\rangle$ を返す。
言語が最初から区別しているので、型が後付けにならない。

\textbf{メリット}：待つか待たないかがコードの見た目に出る。
効果の伝播も、待つ送信（\texttt{now}/\texttt{await}）だけを辿ればよい。
\textbf{デメリット}：三つ書き分ける負担がある。
特に \texttt{future} は \texttt{await} と対で書く必要があり、
書き忘れると「送ったが受け取らない」ことになる。

\subsection{オブジェクト内部の逐次性 ---【採る】}


アクターはメッセージを一つずつ処理する。
フィールドへの代入に排他制御が要らないのはこのためで、
\S\ref{sec:koba} で述べる効果 \texttt{mut} が意味を持つのもこの前提のうえである。

\paragraph{採った理由。}
排他制御を書かずに済むことと、健全性の証明の土台になることの二つである。
\texttt{mut} という効果が意味を持つのも、この前提のうえである ---
「自分の状態を変える」が競合を意味しないのは、一度に一つしか処理しないからである。

\textbf{メリット}：ロックを書かない。データ競合が原理的に起きない。
\textbf{デメリット}：一つのアクターが重い処理をすると、その間そのアクターは止まる。
自分自身への \texttt{now} は原理的に返らない（\S\ref{sec:deadlock}）。

\subsection{\texttt{reply} による返信 ---【採る】}


値は \texttt{return} ではなく \texttt{reply} で返す。
これは ABCL/1 の返信をそのまま受けたものだが、静的型を載せるうえでは厄介な選択である。
\texttt{return} なら関数の型が戻り値型を決めるが、
\texttt{reply} は本体のどこにでも書ける文だからである。
AIPL はメソッドごとに戻り値型を表す型変数 $\rho$ を置き、
本体の各 \texttt{reply(e)} と単一化することで推論する。

% =====================================================================

\paragraph{採った理由。}
源流の考え方であり、型付けの問題も解けたからである。
「メソッドの戻り値型」を \texttt{reply} から推論するために、
メソッドごとに戻り値型変数 $\rho$ を置く。
$\rho$ を\textbf{型スキームに含めない}のが要点で、
含めると \texttt{reply} の型と結べなくなる（\S\ref{sec:koba} の多相メソッドの項）。

\textbf{メリット}：戻り値型の注釈が要らない。書けば強い規律になる。
\textbf{デメリット}：単相に限られる。同じメソッドを違う型で使い回せない。

\subsection{返信先の一級性 ---【採る】}
\label{sec:reply1st}

ABCL/1 では \emph{now} 型送信で暗黙に作られる返信先が値であり、他へ渡せる。
つまり「A が受けた依頼に B が代わりに答える」委譲が書ける。
AIPL もこれを持つ。

ただし、素朴に値にすると「ちょうど一度だけ返信する」が保てない ---
本体に現れる \texttt{reply} を構文で数える方式は、
返信先を渡せるようにした途端に無力になる。
\textbf{線形型}がその答えであった（\S\ref{sec:koba}）。

\paragraph{仕様。}
\texttt{replyto} が、いま処理しているメッセージの返信先を値として返す。
型は \texttt{reply}$\langle\rho\rangle$（$\rho$ はそのメソッドの戻り値型）。
\texttt{answer(r, v)} で返信する（\texttt{reply(v)} はこの糖衣にあたる）。
引数の型注釈 \texttt{reply} で、他のアクターへ渡せる。

\begin{aipl}
class Backend {
  method handle(r: reply, x) : unit !{} {
    answer(r, x * 2);            // Frontend の依頼者へ直接返す
  }
}
class Frontend {
  var b = new Backend();
  method ask(x) : int !{} {
    send b.handle(replyto, x);   // 返信先を渡して義務を移す。自分は reply しない
  }
}
var f = new Frontend();
print("ans = " ++ (now f.ask(21) timeout 500 else -1));
\end{aipl}
\begin{term}
ans = 42
\end{term}

\textbf{委譲が書けて、なお「ちょうど一度」が保たれる。}
これが ABCL/1 に無かったものである。

\paragraph{線形性の規律。}
\texttt{replyto} で義務が生まれ、
\texttt{answer} するか送信の引数に渡すと義務が消える（相手へ移る）。
同じ返信先を二度使ってはならず、
メソッドを抜けるとき義務が残っていてはならず、
\texttt{if} の二つの枝は同じ状態で合流する。

\begin{term}
[Type error] method B19.h returns without answering r
[Type error] reply destination r is used twice
[Type error] the two branches disagree on which reply destinations are owed
\end{term}

\begin{remark}
実装で要点になったのは\textbf{状態の持ち方}である。
「まだ答えていない」集合から消すだけだと、
\textbf{二度渡しがただの変数参照に見えて素通りする}。
$(\text{owed}, \text{spent})$ の対にして初めて捕まった。
線形型は「使ったら消す」ではなく「使ったことを憶える」ことで成り立つ。
\end{remark}

\paragraph{採った理由。}
委譲は実際に要る --- 窓口役が受けて、専門役が直接答える形は、
エージェントでもサービスでも普通に現れる。
そして「ちょうど一度」は線形型で守れる。
落とすべき理由が型の側に無くなった。

\textbf{メリット}：委譲が書けて、なお返信の取りこぼしと二重返信が型で止まる。
\textbf{デメリット}：線形性の検査はメソッド内で閉じている。
返信先を配列に入れて回すような書き方は追えない。

\paragraph{副産物 --- \texttt{self} と \texttt{sender} は値ではない。}
返信先は一級だが、\texttt{self} と \texttt{sender} は
\textbf{送信の宛先としてしか書けない}。

\begin{term}
send self.h(1);          -> OK
send e.take(self, x);    -> PARSE_ERROR
\end{term}

そのため、二者が互いに \texttt{now} で待ち合う循環を、
現在の構文では素直に組み立てられない。
デッドロック自由性はこれとは別に型で保証しているが（\S\ref{sec:dl}）、
\textbf{構文の側でも待ち合いが作りにくい}。

\subsection{\textit{express mode} による割り込み ---【落とす】}


ABCL/1 は ordinary モードとは別の待ち行列を持ち、
express メッセージは実行中の処理を中断して割り込む。

\paragraph{落とした理由。}
\S\ref{sec:abcl} の逐次性が破れる。
「メソッド本体の実行が最後まで途切れない」という前提のうえに、
フィールドの不変条件も、効果の集計も、型健全性の証明も立っている。
割り込みを入れると、本体の途中で状態が観測されうるので、
これらをすべて「割り込み点で成り立つ不変条件」に書き直す必要がある。

\paragraph{落とすと何が書けなくなるか。}
中止・優先処理・監視である。
AIPL では \texttt{select} と期限で近いことを書くが、
\textbf{実行中の処理を止めることはできない}。
できるのは「次に何を受けるか」を選ぶことだけである。

\textbf{メリット（落としたことの）}：逐次性が保たれ、健全性の証明が単純なままである。
\texttt{mut} が競合を意味しないという前提も崩れない。
\textbf{デメリット}：緊急停止のような「今すぐ割り込む」記述ができない。
\texttt{select} と期限で近いことは書けるが、
処理中のメッセージを中断することはできない。

\subsection{\texttt{script} の \texttt{where} ガード ---【落とす】}


ABCL/1 は受信のパターンに \texttt{where} 節を書け、
「在庫があるときだけ購入メッセージを受理する」といった記述ができた。
AIPL の \texttt{select} はパターンをメソッド名と引数の並びに限り、\texttt{where} を持たない。

\paragraph{落とした理由。}
進行定理に枝が一本増えるからである。
どのガードも成立せず受理できない構成が現れうるので、
「終状態」「簡約できる」に加えて「受理できるメッセージが無い」を扱う必要が出る。
さらにガードの純粋性という新しい要求も生じる（ガードが状態を変えてはならない）。

\paragraph{落とすと書き方がどう変わるか。}
受理してから本体で分岐する。

\begin{aipl}
class Shop {
  var stock = 2;
  method buy(n) : string !{mut} {
    if (stock >= n) { stock = stock - n; reply("ok, left=" ++ stock); }
    else { reply("sold out"); }
  }
}
var s = new Shop();
print(now s.buy(1) timeout 200 else "?");
print(now s.buy(5) timeout 200 else "?");
\end{aipl}
\begin{term}
ok, left=1
sold out
\end{term}

違いは\textbf{受理してしまう}点にある。
\texttt{where} なら「受理しない＝送り主は待ち続ける（あるいは他が受理する）」であったのが、
AIPL では「受理して、断りの返信をする」になる。
在庫が戻るまで待たせたい場合、その待ちはプログラマが書くことになる。

% =====================================================================

\textbf{メリット（落としたことの）}：進行定理の枝が増えない。
ガードの純粋性という新しい要求も生じない。
受理条件が本体の \texttt{if} として見えるので、読むときに一箇所で済む。
\textbf{デメリット}：「受理しない」が書けない。
在庫が戻るまで待たせたい場合、その待ちはプログラマが書くことになる。

\subsection{ABCL/1 からの採否 --- まとめ}

\begin{center}
\small
\begin{tabular}{p{0.22\textwidth}p{0.08\textwidth}p{0.30\textwidth}p{0.30\textwidth}}
\toprule
機構 & 採否 & 主な理由 & 代価 \\
\midrule
三種の送信 & 採る & 設計上必要な区別で、型でも別物 & 書き分けの負担 \\
内部の逐次性 & 採る & 排他制御が要らず、証明の土台になる & 一つのアクターは詰まりうる \\
\texttt{reply} & 採る & $\rho$ で型付けできた & 単相のみ \\
返信先の一級性 & 採る & 委譲が要り、線形型で守れる & 検査はメソッド内で閉じる \\
\addlinespace
\textit{express mode} & 落とす & 逐次性が破れ、証明を書き直す必要 & 割り込みが書けない \\
\texttt{where} ガード & 落とす & 進行定理に枝が増え、純粋性が要る & 「受理しない」が書けない \\
\bottomrule
\end{tabular}
\end{center}

% =====================================================================
\section{小林らの型システムの機能と、その採否}
\label{sec:koba}

小林直樹らの並行計算のための型システムの系譜から、
AIPL が何を採り何を落としたかを、同じ形で述べる。

\subsection{資源使用解析（効果）---【採る】}
\label{sec:eff}


メソッドは自分が世界に及ぼす作用を宣言できる。
効果は 8 種類の固定語彙である。

\begin{center}
\begin{tabular}{ll@{\qquad}ll}
\toprule
\texttt{mut}  & 自分の状態を変える    & \texttt{net} & ネットワークに出る \\
\texttt{time} & 時間に依存する        & \texttt{ai}  & 言語モデルを呼ぶ \\
\texttt{io}   & 入出力                & \texttt{fs}  & ファイルを触る \\
\texttt{mem}  & 動的に割り付ける      & \texttt{log} & ログを出す \\
\bottomrule
\end{tabular}
\end{center}

\begin{aipl}
class Store {
  var n = 0;
  method bump()  : int !{mut}  { n = n + 1; reply(n); }   // 状態を変える
  method peek()  : int !{}     { reply(n); }              // 純粋
  method show()  : unit !{log} { print("n=" ++ n); }
}
\end{aipl}

書かなければ本体から推論し、書いた場合は本体から推論した集合と照合する。
足りなければ型エラーになる。

\paragraph{採った理由。}
これが \S\ref{sec:intro} の動機に対する直接の答だからである。
実時間側のアクターに \texttt{!\{time, log\}} と書いておけば、
そこから \texttt{ai} や \texttt{net} を持つメソッドを\textbf{同期で}呼ぶことが型エラーになる。

\begin{term}
[Type error] method Ctrl.viaNow declares {log, time}
             but its body has {ai, log, net} (missing {ai, net})
\end{term}

小林の資源使用解析に比べれば素朴である --- 効果は\textbf{順序を持たない集合}で、
「開いたら閉じる」のような性質は書けない。
そこは今後の課題として残っている（\S\ref{sec:discuss}）。

\paragraph{規律の細部。}
効果を引き継ぐのは\textbf{待つ}呼び出しだけである。
\texttt{send} は送って待たないので引き継がない。
この判断は明示的なもので、逆に \texttt{send} も引き継ぐようにして
手元の \texttt{.aipl} 543 本を通してみたところ、
落ちたのは\textbf{この規律を説明している例題ただ一本}であった。

\paragraph{採った理由。}
実時間側とエージェント側を分ける\textbf{唯一の手段}だからである。
「このメソッドは機外のモデルを呼ぶ」を型に出さない限り、
実時間で動くアクターが知らずにモデルを待つのを止められない。
本格的な資源使用解析（使用の\textbf{順序}まで追う）ではなく、
\textbf{集合}という素朴な形にしたのは、推論が軽く説明しやすいためである。

\textbf{メリット}：注釈が要らない（推論に任せられる）。
\texttt{now}/\texttt{await} で伝播するので、一段隔てて呼んでも隠せない。
メッシュ配備では、そのまま\textbf{ノード境界の入国審査}になる（\S\ref{sec:mesh}）。
\textbf{デメリット}：集合なので\textbf{順序を表せない}。
「取ったら放す」が言えないので、次項を別に足すことになった。

\subsection{効果に順序を入れる ---【採る（拡張して）】}

\label{sec:res}

効果は\textbf{順序を持たない集合}であり、
「開いたら閉じる」「取得したら解放する」が書けなかった。
小林の資源使用解析のいちばん素朴な形を入れた。

\begin{aipl}
class Bus {
  var n = 0;
  method read(addr) : int !{mut} {
    acquire("i2c");
    n = n + 1;
    var v = addr * 2;
    release("i2c");           // これを消すと型エラーになる
    reply(v);
  }
}
\end{aipl}

規律は三つで、いずれもメソッド本体を経路に沿って追って検査する。

\begin{enumerate}
\item メソッドを抜けるとき、持っている資源が残っていてはならない
\item 持っていない資源を \texttt{release} してはならない
\item 二重に \texttt{acquire} してはならない
\end{enumerate}

\texttt{if} の二つの枝は同じ持ち物で合流し、
\texttt{while} の本体は持ち物を変えてはならない。

\begin{term}
[Type error] method Dev4.use returns while still holding i2c
[Type error] resource i2c is acquired twice
[Type error] the two branches disagree on which resources are held (i2c)
\end{term}

\paragraph{限界。}
メソッド内で閉じた検査であり、メソッドをまたぐ受け渡しは見ない。
また資源に\textbf{全体の順序を課していない}ので、
取得順序の食い違いによるデッドロックは防げない（\S\ref{sec:deadlock}）。

\paragraph{採った理由。}
効果を集合にした代価をここで払っている。
資源使用解析のいちばん素朴な形 --- 取得と解放の対 --- だけを、
本体の構文を追って検査する。

\textbf{メリット}：デバイスを掴んだまま抜ける経路が静的に消える。
OS を書く目的に直結する。
\textbf{デメリット}：メソッド内で閉じた検査であり、
メソッドをまたぐ受け渡しは見ない。
資源の名前は文字列リテラルに限る。
なお、対の検査だけでは\textbf{逆順の取得}を止められない ---
そのために全体順序を別に入れた（\S\ref{sec:resorder}）。

\subsection{サイズ型・精製型 ---【粗い代用として採る】}
\label{sec:deadline}


\texttt{now} と \texttt{await} には期限を書ける。

\begin{aipl}
class Slow {
  method fast(x) : int { reply(x); }
  method never(x) : int { while 1 > 0 do { } reply(x); }   // 返ってこない
}
var s = new Slow();
print("fast  = " ++ (now s.fast(42) timeout 200 else -1));
print("never = " ++ (now s.never(7) timeout 200 else -1));
\end{aipl}
\begin{term}
fast  = 42
never = -1
\end{term}

\paragraph{採った理由と、その位置づけ。}
小林のサイズ型・精製型は応答時間の上界を\textbf{型で}述べる。
AIPL の期限はそこまで行かず、実行時に打ち切るだけである。
しかし効果は非常に大きい --- 期限のない待ちを構文から外すと、
待ちが「残り時間を減らす簡約」として進むので、
進行定理から「未解決 future の待ちでブロックする」枝が消える。
系としてデッドロック自由が言える。

代償は明確である。まず停止性は言えない。
そして\textbf{デッドロックが消えるのではなく、有限時間の失敗に変わる}。

\begin{aipl}
class Never  { method wait_forever(x) : int !{} { while 1 > 0 do { } reply(x); } }
class Caller {
  var n = new Never();
  method go(x) : int !{} { reply(now n.wait_forever(x) timeout 200 else -1); }
}
var c = new Caller();
print("go = " ++ (now c.go(5) timeout 900 else -999));
\end{aipl}
\begin{term}
go = -1
\end{term}

詰まらない。しかし正しくもない。$-1$ が返るだけである。

\paragraph{この形で採った理由。}
「この待ちは有限時間で終わる」を型で言うには、
本来サイズ型や精製型で上界を表す必要がある。
AIPL はそれを\textbf{実行時に打ち切る}ことで代用している ---
上界を証明する代わりに、上界を課す。
機外のモデルを呼ぶ以上、待ち時間の上界を静的に知ることは原理的にできない、
という事情もある。

\textbf{メリット}：書ける。証明が要らない。
そして\textbf{失敗が必ず有限時間で起きる}。
\textbf{デメリット}：\textbf{保証ではない}。
「詰まらない」のではなく「詰まっても諦める」である。
またミリ秒という単位は実時間側には粗すぎる。

\subsection{失敗を型に出す ---【サイズ型の代用の穴を塞ぐ】}

\label{sec:result}

期限で代用すると、新しい穴が開く ---
\textbf{期限切れの \texttt{else} の値が正常値と同じ型}になってしまう。
$-1$ が返ることと、正しく計算できたことを型が区別しない。

\paragraph{仕様。}
期限つきの待ちで \texttt{else} を書かなければ、
式の型は \texttt{result}$\langle\tau\rangle$ になる。

\begin{center}
\begin{tabular}{ll}
\toprule
書き方 & 型 \\
\midrule
\texttt{now t.m(a) timeout 100 else 0} & $\tau$（従来どおり） \\
\texttt{now t.m(a) timeout 100}        & \texttt{result}$\langle\tau\rangle$（新） \\
\bottomrule
\end{tabular}
\end{center}

この形は\textbf{従来は構文エラーだった}ので、既存プログラムの意味は変わらない。
取り出しは \texttt{is\_ok} / \texttt{timed\_out} / \texttt{value} で行う。

\begin{aipl}
var a = now s.fast(42) timeout 200;      // result<int>
if (is_ok(a)) { print("fast ok " ++ value(a, -1)); }
else          { print("fast timed out"); }
\end{aipl}

\texttt{result}$\langle\tau\rangle$ は演算子に渡せない。

\begin{term}
[Type error] no overload of + matches (result int, int)
\end{term}

\paragraph{効いていること。}
「この値は本当に計算されたのか」をプログラムが問える。
確かめずに流すことができない。
543 本を通した結果、影響を受けたファイルは 0 本であった。

\paragraph{採った理由。}
精製型を入れずに「この値は本当に計算されたのか」を問えるようにする、
いちばん軽い形だからである。

\textbf{メリット}：時間切れの値が正常値のふりをして下流へ流れることが型で止まる。
AI エージェントでは、これがよく効く（\S\ref{sec:agent}）。
\textbf{デメリット}：\texttt{is\_ok} / \texttt{value} を書く手間が増える。
\texttt{else} を書けば従来どおりなので、規律は\textbf{書き手が選ぶ}ことになる。

\subsection{線形型 ---【採る】}


線形型そのものは採っていないが、
「ちょうど一度だけ使う」という規律は \texttt{reply} に持ち込んでいる。

\begin{itemize}
\item \textbf{線形性}：一つのメッセージに対する \texttt{reply} は高々一度。
\item \textbf{被覆}：戻り値型を宣言したら、すべての実行経路で \texttt{reply} に到達する。
\end{itemize}

\begin{term}
[Type error] method f may reply more than once on some path;
             a method must reply at most once
[Type error] method f is declared to reply with int,
             but some execution path does not reply
\end{term}

この二つの実現は\textbf{構文の走査}である。
返信先を値として渡せるようにした途端、これは無力になる ---
渡された先で何回返されるかは走査では分からない。
そこで返信先そのものには\textbf{線形型}を与えている
（\S\ref{sec:reply1st}）---
\texttt{reply}$\langle\rho\rangle$ を「ちょうど一度使う」義務として扱い、
使ったか（\texttt{spent}）と負っているか（\texttt{owed}）の対で追う。

% =====================================================================

\paragraph{採った理由。}
「ちょうど一度」は返信の要であり、
これが無ければ返信先を一級の値にできない。
線形型があるからこそ、ABCL/1 の委譲を型で支えられる。

\textbf{メリット}：委譲が書けて、なお取りこぼしと二重返信が止まる。
\textbf{デメリット}：全面的な線形型ではない。
線形に扱うのは\textbf{返信先だけ}で、通信路一般や資源には及んでいない
（資源の側は \S\ref{sec:resorder} の順序で別に押さえている）。

\subsection{デッドロック自由な型システム ---【採る】}

通信の依存関係に順序（\textbf{義務レベル}）を入れて循環待ちを型で排除する手法は、
AIPL の課題にそのまま噛み合う。

\paragraph{採った理由。}
期限による代用では\textbf{保証にならない}からである。
「詰まっても諦める」と「詰まらない」は別のことで、
実時間 OS を書く目的には後者が要る。

採るにあたって二つ工夫した。
第一に、レベルは\textbf{書かなくてよい} --- 待ちの辺を見て押し上げる不動点で推論する。
明示注釈は固定点として扱い、動かせなければそこが矛盾である。
第二に、デッドロックを\textbf{一つのものとして扱わなかった}。
調べると経路が四つあり、要る道具が違った（\S\ref{sec:dl}）。

\textbf{メリット}：注釈なしで循環待ちが止まる。
推論したレベルはそのまま\textbf{反例の証人}になる（\texttt{AIOS\_SHOW\_LEVELS=1}）。
ノード境界にも同じ形で課せる（\S\ref{sec:nodelevel}）。
\textbf{デメリット}：宛先がクラスとして分かる範囲でしか効かない。
動的な宛先は見えない。
また粒度が \texttt{Class\#method} なので、
\textbf{同じクラスの別インスタンス}を区別できない。

\subsection{交差型による高階モデル検査 ---【落とす】}


交差型による型付けと $\mu$ 計算のモデル検査の等価性は、
検証器を自作せずに借りるための強力な道具である。

\paragraph{落とした理由。}
型推論が一般に決定不能であり、
AIPL の対象（アクターと通信）に対して必要な表現力を超えている。
加えて、反例が AIPL の言葉で出てこない ---
検査器が返すのは高階再帰スキームの反例であって、
「どのアクターがどのメッセージで詰まるか」ではない。
利用者に見せられない反例は、検証としては使えても道具としては使えない。

\textbf{メリット（落としたことの）}：型検査が決定可能なままである。
エラーが AIPL の言葉（アクター名・メソッド名・行）で出る。
\textbf{デメリット}：借りられたはずの検証力を捨てている。

\subsection{多相メソッド ---【落とす】}


メソッドは単相である。同じメソッドを二つの型で使うことはできない。

\begin{aipl}
class Id { method id(x) : any { reply(x); } }
var i = new Id();
print(now i.id(1) timeout 100 else 0);
print(now i.id("s") timeout 100 else "?");
\end{aipl}
\begin{term}
[Type error] line 3, col 7: type mismatch
\end{term}

\paragraph{落とした理由。}
戻り値型を表す $\rho$ を型スキームに含めると、
\textbf{呼び出しごとに $\rho$ が新しくなる}。
すると「このメソッドはどの型を返すか」という制約が呼び出し側へ伝わらず、
\texttt{reply} の地点と結ばれなくなる。
$\rho$ はクラスのメソッド表に属する量であって、多相化の対象ではない。
単相のままにしておくのが、\texttt{reply} からの推論と型健全性の証明の双方に対して整合的である。

\paragraph{代価。}
上の \texttt{Id} のような汎用のアクターが書けない。
実際には \texttt{any} を使って型検査を素通りさせることになるが、
それは検査を放棄しているだけである。

\subsection{完全な所有権の型システム ---【落とす】}


Rust 相当の借用検査は作らない。
CHC への還元に必要な範囲（別名が無いこと）は、
アクターがメッセージでしか相互作用しないという性質から既に成立しているためである。

\textbf{メリット（落としたことの）}：借用の注釈を書かずに済む。
\textbf{デメリット}：資源の別名管理が要る場面 ---
たとえば\textbf{資源がアクターの実体}である場合 --- には届かない
（\S\ref{sec:discuss}）。

\subsection{小林の系譜からの採否 --- まとめ}

\begin{center}
\small
\begin{tabular}{p{0.22\textwidth}p{0.10\textwidth}p{0.28\textwidth}p{0.30\textwidth}}
\toprule
機構 & 採否 & 主な理由 & 代価 \\
\midrule
資源使用解析（効果） & 採る & 実時間側とエージェント側を分ける唯一の手段 & 集合なので順序を表せない \\
効果の順序 & 採る & 集合にした代価を払う & メソッド内で閉じる \\
サイズ型・精製型 & 代用 & 上界を証明する代わりに課す & 保証ではない \\
線形型 & 採る & 「ちょうど一度」は返信の要 & 返信先だけに限る \\
デッドロック自由な型 & 採る & 代用では保証にならない & 宛先が静的に分かる範囲 \\
\addlinespace
交差型モデル検査 & 落とす & 推論が決定不能。反例が言語の言葉で出ない & 検証力を捨てている \\
多相メソッド & 落とす & $\rho$ をスキームに含めると \texttt{reply} と結べない & 汎用アクターが書けない \\
完全な所有権型 & 落とす & 別名の無さはアクターの性質から既に成立 & 実体ごとの管理に届かない \\
\bottomrule
\end{tabular}
\end{center}

% =====================================================================
\section{どちらの源流にも無いもの}
\label{sec:orig}

三つある。いずれも AIPL の目的 --- 分散 AI エージェントと実時間 OS ---
から要請されたものである。

\subsection{セッション型 ---【採る】}

\label{sec:session}

経路 3 には、もう一つの側面がある --- \textbf{順序}である。
「開いてから読み、最後に閉じる」という約束は、
期限でも送り手の存在でも表せない。

\paragraph{既にあったものを持ち上げた。}
AIPL には\textbf{実行時}のセッションプロトコルが既にあった ---
\texttt{protocol\_define} / \texttt{protocol\_start} / \texttt{protocol\_end} で、
Coq でも \texttt{send\_ok\_unique}、\texttt{send\_ok\_advances\_one}、
\texttt{wrong\_label\_rejected} などが検証されている。
新しく作ったのではなく、\textbf{同じ宣言を型検査の側が読む}ようにした。

\begin{aipl}
protocol_define("Pipeline", "f.get -> r.scale -> st.put");
var sid = protocol_start("Pipeline");

var a = now f.get(4)    timeout 200 else 0;
var b = now r.scale(a)  timeout 200 else 0;
var c = now st.put(b)   timeout 200 else 0;
protocol_end(sid);
\end{aipl}
\begin{term}
pipeline = 41
\end{term}

順序を取り違えると、走らせる前に出る。

\begin{term}
[Type error] protocol P29: expected b.step2 but the program sends c.step3 here
\end{term}

\paragraph{採った理由。}
「開いてから読み、最後に閉じる」という約束は、
期限でも送り手の存在でも表せない。
そして AIPL には\textbf{実行時の}セッションプロトコルが既にあった ---
新しい宣言を作る必要が無かった。

\textbf{メリット}：手順の取り違えが実行前に出る。
実行時の検査と\textbf{同じ宣言}を使うので、二重に書く必要が無い。
\textbf{デメリット}：手順は\textbf{名前}で役を指すので、
名前が静的に解決できる範囲でしか効かない（次項）。

\subsubsection*{アクターをまたぐセッション --- 振る舞い展開}

\label{sec:crosssession}

実際のプログラムでは、手順は一つの並びに収まらない。
セッションを使う 13 本のうち\textbf{6 本}がアクターをまたいでいた。
そして 6 本とも、全部\textbf{同じ形}をしていた。

\begin{aipl}
protocol_define("FutureTimelineJa",
  "main_thread.run -> planner.plan -> solver.solve
   -> reviewer.brief -> reviewer.review");
var result = aios_now("main", "run", problem);      // 手順1だけがここ
\end{aipl}

残りの手順は \texttt{main\_thread.run} の本体にある。

\begin{aipl}
var plan          = aios_now("planner", "plan", problem);          // 手順2
var answer_future = aios_future("solver", "solve", problem, plan);  // 手順3
var brief_future  = aios_future("reviewer", "brief", problem, plan);// 手順4
var verdict       = aios_now("reviewer", "review", ...);            // 手順5
\end{aipl}

\textbf{プログラム順にぴったり並んでいる}。
トップが窓口役をひとつ呼び、窓口役の本体が残りを順に送る ---
6 本ともこの形であった。

\paragraph{同期の呼び出しは、その場に差し込める。}
必要だったのは新しい型ではなく、\textbf{展開}であった。
\texttt{now} / \texttt{aios\_now} は、呼び先が呼び出し側の続きより
\textbf{先に走り切る}。だから呼び先の送信列を、呼び出しの位置に
そのまま差し込んでよい。これで、またいだセッションも一本の並びになる。

非同期（\texttt{send} / \texttt{future}）は差し込まない ---
いつ走るか決まらないからである。
呼び先が手順を含んでいたら「静的に並べられない」と印を付け、
完了の判定を実行時に任せる。
送信そのものは呼んだ位置で起きるので、
\textbf{順序の検査には非同期でも支障がない} ---
上の \texttt{aios\_future} 二つが手順 3・4 として正しく並ぶのは、
このためである。

\paragraph{名前を解決する。}
手順は \texttt{"main\_thread.run"} とアクターの変数名で書かれ、
送信は \texttt{aios\_now("main", "run", ...)} とサービス名で書かれる。
両者は \texttt{aios\_register\_service("main", "main\_thread")} で
結ばれている --- 前節で「綴りの違い」として避けた食い違いは、
\textbf{解決表を作れば消える}ものであった。
クラスも同じで、\texttt{var planner = new PlannerActor();} と
引数の型注釈 \texttt{method run(f: Fetch, ...)} からたどれる。

\paragraph{結果。}
6 本とも、\textbf{手順が最後まで静的に並んだ}。
窓口役の本体で順序を入れ替えると、走らせる前に出る。

\begin{term}
[Type error] protocol ThreeActorSolve: expected solver.solve
             but the program sends reviewer.review here
\end{term}

手順をひとつ落とした場合も、今度は完了検査が働く。

\begin{term}
[Type error] protocol ThreeActorSolve is incomplete at protocol_end;
             next expected reviewer.review
\end{term}

582 本を旧版の型検査器と\textbf{一件ずつ}突き合わせた結果、
判定が変わったファイルは\textbf{0 本}であった ---
新しく捕まえるものだけが増え、通っていたものは通り続けている。

\paragraph{実装が教えてくれた限界。}
最初に書いたサンプルは、窓口役が \texttt{Fetch} と \texttt{Refine} を
\textbf{フィールドとして}持つ形であった。型検査は通ったが、走らせると止まった。

\begin{term}
protocol Cross violation: expected f.get, got c#f.get
\end{term}

実行時はフィールドのアクターを\textbf{持ち主で修飾した名前}（\texttt{c\#f}）
で呼ぶ。静的な側はフィールド名 \texttt{f} しか知らない。
クラスに実体はいくつでもありうるので、静的な側が \texttt{c\#f} を
作ることは原理的にできない。
\textbf{手順が名前で役を指す以上、フィールドに隠れたアクターは役になれない}。
サンプルは、アクターを引数で渡す形（\texttt{run(f: Fetch, r: Refine, k)}）に
書き換えた --- こちらは静的にも実行時にも同じ名前で呼ばれる。

\begin{finding}
「アクターをまたぐから静的には追えない」は、\textbf{型の限界ではなく
検査の範囲の狭さ}だった。
必要だったのは多者間セッション型でも委譲でもなく、
\textbf{呼び先の送信列を差し込むこと}と\textbf{名前の解決表}である。
プログラムは何も変えていない。

そして限界は、予想していたところ（またぐこと）ではなく
\textbf{名前}に現れた ---
実行時と静的な側で、同じアクターの呼び名が違っていた。
\end{finding}

\subsection{資源への全体順序 ---【採る】}

\label{sec:resorder}

デッドロックには、もう一本の経路がある。
これは待ちの閉路でも、返されない返信でもない ---
\textbf{二つのアクターが、同じ二つの資源を逆の順序で取る}という形である。

\S\ref{sec:res} の対の検査は、この形を素通しする。
どちらのメソッドも、取ったものは全部返しているからである。

\begin{aipl}
class Writer {
  method run() : int !{mut} {
    acquire("db"); acquire("cache");
    release("cache"); release("db"); reply(1);
  }
}
class Reader {
  method run() : int !{mut} {
    acquire("cache"); acquire("db");     // 逆順
    release("db"); release("cache"); reply(2);
  }
}
\end{aipl}

\paragraph{入れ子が順序を語っている。}
新しい注釈は要らなかった。
\textbf{$r$ を持っているあいだに $s$ を取った}という事実が、
そのまま $r < s$ という辺である。
プログラム全体からこの辺を集め、\textbf{閉路が無いこと}を見る。

\begin{term}
[Type error] resource order: cache -> db -> cache is circular;
  cache before db (in Reader.run); db before cache (in Writer.run).
  Acquiring in opposite orders deadlocks
\end{term}

閉路の各辺が\textbf{どこで生まれたか}を持っておくと、
「逆順に取っている二か所」がそのまま出る。
デッドロックの報告としては、これがいちばん役に立つ形である。

\paragraph{閉路が無い＝全体順序が作れる。}
これは \S\ref{sec:dl3} の義務レベルと同じ構造をしている。
閉路が無いなら位相順序が取れ、その高さがそのまま
「どれを先に取るか」の番号になる。
レベルはここでも\textbf{証人}である（\texttt{AIOS\_SHOW\_LEVELS=1}）。

\begin{term}
[resource] bus                  @0
[resource] dev                  @1
\end{term}

\paragraph{書きたい人のために。}
推論に任せず \texttt{resource\_order("bus -> dev")} と書くこともできる。
宣言は辺として同じ表に入るだけなので、
宣言と実際の取得が食い違えば、それも閉路として出る ---
違反を見るための別の仕掛けは要らなかった。

\paragraph{見えないところ。}
この検査が追えるのは、資源の名前が\textbf{文字列リテラル}である
\texttt{acquire} だけである。
変数に入った名前は静的には分からない
（\texttt{AIOS\_STRICT\_RESOURCE=1} でその場所を知らせる）。
そこで、宣言された順序は\textbf{実行時にも効かせる}ことにした ---
下位を持ったまま上位へ、の向きだけを許す。

\begin{term}
acquire: dev is held, so bus must not be acquired now (declared order)
\end{term}

\begin{finding}
資源の順序は、\textbf{書かせるものではなく読み取るもの}だった。
入れ子はプログラムの中に既にあり、
それを集めれば順序の主張になる。
そして閉路の検査は、\textbf{反例をそのまま報告できる}という副産物を持つ ---
「$a$ の次に $b$ を取っている場所」と「$b$ の次に $a$ を取っている場所」の
両方が、辺の出どころとして残っているからである。

582 本の \texttt{.aipl} に対して、この検査で新たに落ちるものは
\textbf{0 本}であった。
\end{finding}

\subsection{メッシュ配備 ---【採る】}
\label{sec:mesh}


実行基盤として複数の Xinu ノードをメッシュで繋ぐ構想がある。
そこではアクターのソースをネットワーク越しに送り、
相手先で組み立てて動かす記述が要る。

\begin{aipl}
node_allow("rt0",   "mut, log, time");        // 実時間ノードが許す効果
node_allow("edge0", "mut, log, ai, net");

var h = deploy("rt0", "Servo", "servo1");     // ソースを送り、相手先で JIT
send remote("rt0", "servo1").step(3);
\end{aipl}
\begin{term}
[deploy] Servo -> rt0/servo1 (compiled at destination)
pos=3
[Top-level CallStmt error] deploy: node rt0 does not accept effect(s) {ai, net}
                           required by Planner
\end{term}

要点は、\textbf{効果注釈がコードと一緒に運ばれ、受け入れ側の方針と照合される}ことである。
実時間ノードに \texttt{ai} や \texttt{net} を持つコードを配ろうとすると、
配備そのものが失敗する。
\S\ref{sec:koba} の効果が、単一ノード内の規律から\textbf{ノード間の入国審査}になる。

\paragraph{送る荷物。}
クラスを抽象構文木から組み立て直すのではなく、\textbf{ソースの原文}を送る。
組み立て直すと \texttt{select} の中身が落ちるためである。
輸送は現状ひとつのプロセス内で模しているが、
言語から見える形 --- 何を送り、相手が何を検査するか --- は実機と同じである。

% =====================================================================

\paragraph{採った理由。}
複数の Xinu ノードでの実行に要るからである。
そして\textbf{効果注釈がそのまま入国審査になる} ---
実時間ノードに \texttt{ai} や \texttt{net} を持つコードは配れない。
言語の側に既にあるものが、そのまま分散の規律になった。

\textbf{メリット}：配るものがソースなので、受け入れ側が\textbf{自分の方針で}検査できる。
効果とレベルがコードと一緒に運ばれる。
\textbf{デメリット}：受け入れ側に処理系が要る。
また輸送は現状ひとつのプロセス内で模しており、実機のノード間転送はこれからである。

% =====================================================================
\section{採否の対応表}
\label{sec:table}

\begin{center}
\small
\begin{longtable}{p{0.20\textwidth}p{0.07\textwidth}p{0.30\textwidth}p{0.33\textwidth}}
\toprule
機構 & 採否 & 理由 & AIPL での姿 \\
\midrule
\endhead
\multicolumn{4}{l}{\textbf{ABCL/1 から}}\\
\midrule
三種の送信 & 採る & 「待たない／待つ／後で」は設計上必要な区別 &
  \texttt{send} / \texttt{now} / \texttt{future}+\texttt{await} \\
内部の逐次性 & 採る & 排他制御を書かずに済む。証明の土台 & 一度に一つ処理 \\
\texttt{reply} による返信 & 採る & 源流の考え方。型付けは $\rho$ で解決 &
  \texttt{reply(e)}、戻り値型は推論または注釈 \\
返信先の一級性 & 採る & 委譲は要る。線形型で「ちょうど一度」を守れる &
  \texttt{replyto} / \texttt{answer} / 引数型 \texttt{reply} \\
\addlinespace
\textit{express mode} & 落とす & 逐次性が破れ、証明を全面的に書き直す必要 &
  無し。\texttt{select}＋期限で近いことのみ \\
\texttt{where} ガード & 落とす & 進行定理に枝が増え、ガードの純粋性が要る &
  \texttt{select} は名前と引数のみ。本体で分岐 \\
\midrule
\multicolumn{4}{l}{\textbf{小林らの型システムから}}\\
\midrule
資源使用解析 & 採る & 実時間側とエージェント側を分ける唯一の手段 &
  8 種の効果集合 \\
効果の順序 & 採る & 集合では「取ったら放す」を表せない &
  \texttt{acquire} / \texttt{release} の対＋全体順序 \\
サイズ型・精製型 & 代用 & 上界を証明する代わりに実行時に打ち切る &
  \texttt{timeout $n$ else $e$}、\texttt{result}$\langle\tau\rangle$ \\
線形型 & 採る & 「ちょうど一度」は返信の要 &
  返信先の線形性検査（\texttt{owed}, \texttt{spent}） \\
デッドロック自由な型 & 採る & 期限による代用は保証にならない &
  義務レベル \texttt{@$n$}（推論。ノード境界にも課す） \\
\addlinespace
交差型モデル検査 & 落とす & 推論が決定不能。反例が言語の言葉で出ない & 無し \\
多相メソッド & 落とす & $\rho$ をスキームに含めると \texttt{reply} と結べない & 単相のみ \\
完全な所有権型 & 落とす & 別名の無さは既に成立 & 無し \\
\midrule
\multicolumn{4}{l}{\textbf{どちらの源流にも無いもの}}\\
\midrule
セッション型 & 採る & 順序の約束は期限でも存在検査でも表せない &
  \texttt{protocol\_define} / \texttt{\_start} / \texttt{\_end}（静的にも読む） \\
資源への全体順序 & 採る & 対の検査では逆順の取得を止められない &
  取得の入れ子から辺を集めて閉路を見る \\
メッシュ配備 & 採る & 複数 Xinu ノードでの実行に要る。
  効果注釈が入国審査になる &
  \texttt{source\_of} / \texttt{node\_allow} / \texttt{deploy} \\
\bottomrule
\end{longtable}
\end{center}

\begin{finding}
表を通して見ると、採否の分かれ目は一つに見える ---
\textbf{反例が AIPL の言葉で出るか}である。
交差型モデル検査を落としたのは検証力が足りないからではなく、
反例が高階再帰スキームの言葉で返るからであった。
逆にデッドロック自由な型を採れたのは、
推論したレベルがそのまま\textbf{「どのメソッドが何を待って詰まるか」}
という形の証人になったからである。
型システムの価値は、通す力ではなく\textbf{落とし方}にある。
\end{finding}

% =====================================================================
\section{デッドロック --- 四つの経路と、それぞれの道具}
\label{sec:dl}

エージェントが止まる、というのはひとつの現象に見える。
実際には、\textbf{ひとつのものとして扱うと答えを誤る}。
待ちが返らなくなる形は四つあり、要る道具が違う。

\subsection{循環待ちの静的検査}
\label{sec:deadlock}


まず、いちばん安く手に入るものから見る。
効果の伝播に使っている辺（待つ側 $\to$ 待たれる側）が、
\textbf{そのまま待ちグラフ}であった。閉路検出はほぼ費用ゼロである。

\begin{term}
[Type error] circular wait: Peer2#bounce;
             a now/await cycle deadlocks unless a deadline breaks it
\end{term}

543 本すべてで閉路が 0 件だったので、既定をエラーにした
（\texttt{AIOS\_LAX\_WAIT=1} で警告に降格）。

\paragraph{ただし、これだけではデッドロック自由にならない。}
次の 10 行は\textbf{型検査を通り}、実行すると詰まる。

\begin{aipl}
class Peer {
  method tick(n) : int !{} { reply(n); }
  method bounce(p: Peer, n) : int !{} {
    if (n > 0) { reply(now p.bounce(p, n - 1) timeout 300 else -1); }
    else { reply(0); }
  }
}
var p = new Peer();
print("bounce = " ++ (now p.bounce(p, 3) timeout 1500 else -999));
\end{aipl}
\begin{term}
bounce = -1
\end{term}

アクターは一度に一つしかメッセージを処理しないので、
自分への \texttt{now} は原理的に返らない。
$-1$ は\textbf{期限が救っただけ}であって、型は何も止めていない
（この例は閉路検査が拾うが、検査はメソッド単位で保守的かつ不完全であり、
遠隔配備先は見えない）。

\subsection{四つの経路}
\label{sec:dl3}


AIPL で待ちが返らなくなる経路は\textbf{四つ}ある。
うち三つは待ちの形をしているが、四つ目はしていない。

\paragraph{経路 1 --- 循環待ち。}
上の \texttt{Peer\#bounce} である。\texttt{now} の輪。

\paragraph{経路 2 --- 返信されない待ち。}
\textbf{閉路が無くても詰まる。}これは実測で見つけた。

\begin{aipl}
class Silent2 { method f(x) : unit !{log} { print("ignored"); } }
class Caller4 {
  var s = new Silent2();
  method go(x) : unit !{log} {
    var r = now s.f(x) timeout 200;      // result<unit>
    print("after " ++ is_ok(r));
  }
}
\end{aipl}
\begin{term}
ignored
after false
\end{term}

待ちグラフは非巡回である。それでも返信は永久に来ない。
\texttt{is\_ok} が \texttt{false} なのは\textbf{期限が救っただけ}である。

\paragraph{経路 3 --- \texttt{select} が来ないメッセージを待つ。}
\texttt{select \{ case j(k) -> ... \}} を期限なしで書くと型検査は通る。
これも閉路ではない。

\paragraph{経路 4 --- 資源を逆の順序で取る。}
待ちの形はしていない。
二つのアクターが同じ二つの資源を逆順に取ると、
どちらのメソッドも「取ったら返す」を守っているのに、
実行すると互いの手放しを待つ。
\S\ref{sec:res} の対の検査は、この形を素通しする。

\paragraph{そして、順序そのもの。}
四つとは別に、\textbf{やり取りの順序}という軸がある。
「開いてから読み、最後に閉じる」という約束は、
待ちの構造でも資源の取得でもない。
これはセッション型で扱う（\S\ref{sec:session}）。

\subsection{四つとも型で押さえる}


\begin{center}
\begin{tabular}{lll}
\toprule
経路 & 型で防げるか & 要る道具 \\
\midrule
1. 循環待ち & \textbf{はい} & 義務レベル（小林 1997/2000/2006） \\
2. 返信されない待ち & \textbf{はい} & \texttt{reply} の全域性 \\
3. \texttt{select} の待ち & 部分的に & 期限の義務づけ＋送り手の存在検査 \\
3'. やり取りの順序 & \textbf{はい} & セッション型（実行時の宣言を静的に検査） \\
4. 取得順序の食い違い & \textbf{はい} & 資源への全体順序（入れ子から辺を集めて閉路を見る） \\
\bottomrule
\end{tabular}
\end{center}

\paragraph{経路 2 --- \texttt{reply} の全域性。}
戻り値型を宣言したメソッドには全経路での \texttt{reply} を課していたが、
\texttt{unit} を返すメソッドが抜け穴になっていた。
\textbf{\texttt{now}/\texttt{await} で待たれるメソッドに限って}、
戻り値型に関わらず課すことにした。
待たれるメソッドの集合は、効果の伝播に使う辺からそのまま取れる。

\begin{term}
[Type error] method Silent2.f is waited on by now/await
             but does not reply on every path
\end{term}

563 本を通した結果、\textbf{影響を受けたファイルは 0 本}であった。

\paragraph{経路 1 --- 義務レベル。}
各メソッドにレベルを付け、\texttt{now}/\texttt{await} は
\textbf{厳密に大きいレベルにしか向かえない}とする。
全メソッドに付けば、待ちのグラフは構成的に非巡回になる。

\begin{aipl}
class Source9 { method get(i) : int @ 3 !{} { reply(i * 10); } }
class Scale9  {
  var s = new Source9();
  method at(i) : int @ 2 !{} { reply((now s.get(i) timeout 100 else 0) + 1); }
}
class Sink9 {
  var c = new Scale9();
  method run(n) : int @ 1 !{log} { ... now c.at(n) ... }
}
\end{aipl}

待ちは $1 \to 2 \to 3$ と一方向に流れる。逆流も同レベルも弾かれる。

\begin{term}
[Type error] obligation level: B9#g (@2) waits on A9#f (@2);
             a wait must go to a strictly higher level
[Type error] obligation level: B10#g (@5) waits on A10#f (@1); ...
\end{term}

\paragraph{注釈は推論できる。}
[25] の助言に従い明示宣言から始めたが、
そこには弱点があった --- \textbf{片方にしか注釈が無い辺は検査されない}。
注釈を付け忘れた組が素通りする。効果注釈と同じ弱点である。

そこで、注釈が無いメソッドには\textbf{レベルを推論する}ようにした。
規則は一つで、待ちの辺 $(a, b)$ について
$\mathrm{level}(b) \le \mathrm{level}(a)$ なら $b$ を押し上げる、を不動点まで繰り返す。
明示注釈は固定点であり、\textbf{押し上げが要るのに動かせなければ、そこが矛盾}である。

\begin{term}
$ AIOS_SHOW_LEVELS=1 tc s10_pipeline.aipl
[level] Sink#run                     @0
[level] Scale#at                     @1
[level] Source#get                   @2
OK
\end{term}

注釈をひとつも書いていないプログラムに、待ちの構造がそのまま出てくる。
混在しても誤検出しない --- \texttt{Mid\#at @5 (declared)} に対して
\texttt{Src\#get} は \texttt{@6} と推論される。
矛盾する注釈は捕まる。

\begin{term}
[Type error] obligation level: Mid2#at (@9) waits on Src2#get (@1);
             a wait must go to a strictly higher level
\end{term}

閉路があると不動点に達しないので、閉路検査を先に走らせて具体的な診断を出す。
これも既存コードへの影響は 0 本である。

\begin{finding}
推論を入れて分かったことがある。
\textbf{レベル推論は、閉路検査に「証拠」を与えたものである}。
閉路が無いことと、レベルの割り当てが存在することは同じである。
違うのは、閉路検査が「無い」としか言わないのに対し、
推論は\textbf{どういう順序で待っているか}を数として見せることである。
そして混在を許すことで、\textbf{人が書いた規律}（明示注釈）と
\textbf{コードが持つ構造}（推論）を突き合わせられるようになった。
\end{finding}

\subsection{「保証できる」と言うための条件}


それでも、\textbf{全言語について}デッドロック自由が言えるわけではない。
条件を明記しておく。

\begin{itemize}
\item \textbf{引数で受け取ったアクターは、注釈があれば覆われている}。
      \texttt{method g(s: Svc1, x)} と書けば
      \texttt{Use1\#g @0 $\to$ Svc1\#f @1} と辺が張られる。
      注釈が無ければ、そもそも「\texttt{target is not actor}」で弾かれる。
      つまりレベルを\textbf{アクター型に載せる必要は無かった}。
\item \textbf{精度が落ちる}。レベルはメソッド単位なので、
      同じクラスの別インスタンス同士（親が子を待つ木構造など）も一律に弾かれる。
      緩めるにはレベル変数が要る。
\item \textbf{ノードをまたぐ待ちは、ノード単位で階層化する}（下記）。
\item \textbf{停止性・飢餓は別問題}。
      \texttt{while 1>0 do \{\}} がメールボックスを塞ぐのはレベルでは止まらない。
\item 経路 3（\texttt{select}）は\textbf{部分的}である。
      送信が動的（\texttt{sender} 経由・実行時に決まる宛先）だと見えない。
      完全には、セッション型が要る。
\end{itemize}

\subsection{経路 3 --- \texttt{select} の待ち}
\label{sec:select3}


三つ目は\textbf{閉路でも返信漏れでもない}。誰も送らないだけで詰まる。
セッション型がなければ完全には扱えないが、
実用上ほとんどを占める二つの形は押さえられる。

\paragraph{規律 1 --- 期限を書く。}
期限の無い \texttt{select} は永久に待ちうる。
\texttt{now}/\texttt{await} と同じ扱いにした --- 既定は警告、
\texttt{AIOS\_STRICT\_DEADLINE=1} でエラーに昇格する。
永久の待ちが有限時間の失敗に変わる。

\paragraph{規律 2 --- 送り手が居るか。}
\texttt{select} の \texttt{case} が待つメッセージを、
プログラムの中で誰かが送っているかを見る。
送り手がどこにも無ければ、その \texttt{case} は永久に発火しない。

\begin{term}
[warn] select waits for W2.ghost but nothing in this program sends it
\end{term}

\paragraph{外部の送り手。}
この検査は最初、563 本で 2 件の警告を出した。
調べると\textbf{どちらも \texttt{web\_expose} で外へ公開している例題}で、
送り手は HTTP の向こう側にいた。
\texttt{web\_expose} / \texttt{web\_listen} / \texttt{deploy} / \texttt{remote} を
使うプログラムでは、この検査は当てにならないので行わない。
結果、563 本での警告は\textbf{0 件}になった。

\begin{finding}
「誰も送らない」を静的に言うには、\textbf{プログラムの外側を数えなければならない}。
効果や返信は言語の中で閉じているが、
メッセージの送り手は Web の向こうにも、別のノードにもいる。
この検査が既定で警告なのはそのためである ---
\textbf{言語が世界の全部を見ているわけではない}という事実が、
そのまま検査の強さの上限になる。
\end{finding}

\subsection{ノードをまたぐ待ち}
\label{sec:nodelevel}

レベルを「アクター型に載せる」ことも考えられるが、
実際に穴を数えてみると\textbf{そこではなかった}。
クラス注釈つきの引数は既に辺を張っており、
注釈が無ければそもそも弾かれる。
残っていた穴は\textbf{ノードをまたぐ待ち}だけであった。

宛先の実体は別ノードにあり、静的には見えない。
そこで\textbf{ノード単位で階層化}する。

\begin{aipl}
class Servo {
  var pos = 0;
  method step(d) : int @ 12 !{mut, log} { pos = pos + d; print("pos=" ++ pos); reply(pos); }
}
class Controller {
  method drive(d) : int @ 2 !{net} {          // 上位は低いレベル
    reply(now remote("rt0","servo1").step(d) timeout 200 else -1);
  }
}
node_level("rt0", 10);              // 実時間ノードは葉（高いレベル）
node_allow("rt0", "mut, log, time");
deploy("rt0", "Servo", "servo1");
\end{aipl}
\begin{term}
[deploy] Servo -> rt0/servo1 (compiled at destination)
pos=3
drive = 3
\end{term}

そのノードへの待ちは必ず上へ向かうことを要求する。

\begin{term}
[Type error] obligation level: Ctl2#go (@7) waits on node rt0 (floor @1);
             a wait across nodes must go up
\end{term}

配備の側でも照合する --- 下限より低いレベルのコードを置くと、
そのノードへの待ちが上へ向かうという約束が破れる。

\begin{term}
deploy: Servo.step is @3 but node rt0 has floor @10
\end{term}

\textbf{効果を \texttt{node\_allow} で国境で照合しているのと、まったく同じ形}である。
効果もレベルも、コードと一緒に運ばれ、受け入れ側の方針と突き合わされる。

\begin{remark}
併せて一つ直した。\texttt{any} がどの型とも単一化しなかったため、
\texttt{now remote(n,a).m(x) timeout 100 else 0} が
\texttt{type mismatch} になり\textbf{そもそも書けなかった}。
\texttt{any} は「型が分からない」の印であって、合わない印ではない。
563 本への影響は 0 本であった。
\end{remark}

\paragraph{到達点。}
\textbf{待ちの宛先がクラスとして分かるか、宛先ノードのレベルが宣言されている範囲では、
デッドロック自由が型で保証される。}
注釈は要らない --- 推論が埋める。
注釈を書けば、コードが持つ構造より\textbf{強い規律}を課せる。
その外側は、期限が有限時間の失敗に変換する。
資源の側は \S\ref{sec:resorder} で全体順序を課したので、
取得順序のデッドロックも同じ形（辺を集めて閉路を見る）で止まる。

「アクター型にレベルを載せる」ことも検討したが、\textbf{要らなかった} ---
型注釈つきの引数が既に辺を作っているので、
足りなかったのはノードをまたぐ待ちだけであった。

% =====================================================================

% =====================================================================
\section{型付き AI エージェントを書く}
\label{sec:agent}

ここまでの機構は、それぞれ別の動機から入れたものである。
本節では、それらを\textbf{一つのプログラムの上で}働かせてみる ---
立案・求解・査読の三役からなる、ごく普通の AI エージェントである。

\subsection{三役のパイプライン}

\begin{aipl}
class Planner {
  method plan(q) : string !{ai, net} {
    reply(ai_call_with_system("You are a planner. Output one line.", q));
  }
}
class Solver {
  method solve(p) : string !{ai, net} {
    reply(ai_call_with_system("You are a solver. Answer concisely.", p));
  }
}
class Reviewer {
  method review(a) : string !{ai, net} {
    reply(ai_call_with_system("You are a reviewer. Reply OK or NG only.", a));
  }
}
var planner  = new Planner();
var solver   = new Solver();
var reviewer = new Reviewer();

protocol_define("Agent", "planner.plan -> solver.solve -> reviewer.review");
var sid = protocol_start("Agent");

var p = now planner.plan("What is 12 * 11?") timeout 20000;   // result<string>
var a = now solver.solve(value(p, "")) timeout 20000;
var v = now reviewer.review(value(a, "")) timeout 20000;

protocol_end(sid);
print("verdict = " ++ value(v, "(timed out)"));
\end{aipl}

型の付いていない書き方と較べて増えているのは、次の四つだけである。

\begin{center}
\begin{tabular}{lll}
\toprule
書いたもの & 何を言っているか & 節 \\
\midrule
\texttt{!\{ai, net\}} & このメソッドは機外のモデルを呼ぶ & \S\ref{sec:eff} \\
\texttt{timeout 20000} & 20 秒待って、返らなければ諦める & \S\ref{sec:deadline} \\
\texttt{result}$\langle$\texttt{string}$\rangle$（\texttt{else} 無し） & 諦めた場合と成功した場合を型で分ける & \S\ref{sec:result} \\
\texttt{protocol\_define} & 立案 $\to$ 求解 $\to$ 査読の順に進む & \S\ref{sec:session} \\
\bottomrule
\end{tabular}
\end{center}

\subsection{何が止まるか --- 四つとも壊してみた}

「型が付いている」と言うだけでは意味が無い。
上のプログラムを四通りに壊し、実際に処理系へ通した。

\paragraph{(1) 期限を外す。}
\texttt{timeout 20000} を消すと、まず警告が出る。

\begin{term}
[warn] now without a deadline waits forever;
       write `now ... timeout <ms> else <expr>`
[Type error] no overload of value matches (string, string)
\end{term}

面白いのは\textbf{その下の型エラー}である。
期限を消すと待ちの型が \texttt{result}$\langle$\texttt{string}$\rangle$ ではなく
\texttt{string} になるので、後続の \texttt{value(p, "")} が型として合わなくなる。
期限は\textbf{飾りではなく、後ろの型を支えている}。

\paragraph{(2) \texttt{ai} 効果を隠す。}
\texttt{!\{ai, net\}} を \texttt{!\{net\}} に書き換える。

\begin{term}
[Type error] method Planner.plan declares {net} but its body has {ai, net}
             (missing {ai})
\end{term}

これは \S\ref{sec:intro} の動機そのものである ---
実時間側のアクターが「これはモデルを呼びません」と偽ることができない。
効果は \texttt{now}/\texttt{await} で呼び先から伝播するので、
\textbf{一段隔てて呼んでも隠せない}。

\paragraph{(3) 役割の順序を入れ替える。}
査読を求解より先に呼ぶ。

\begin{term}
[Type error] protocol Agent: expected solver.solve
             but the program sends reviewer.review here
\end{term}

\paragraph{(4) 失敗を確かめずに次へ渡す。}
\texttt{value(p, "")} を \texttt{p} に書き換える ---
つまり \texttt{result}$\langle$\texttt{string}$\rangle$ をそのまま次の役割へ流す。

\begin{term}
[Type error] type mismatch
\end{term}

時間切れの値が正常な答えのふりをして下流へ流れる、という
エージェント特有の壊れ方が、ここで止まる。

\begin{finding}
四つのうち三つは\textbf{別々の動機で入れた機構}である ---
効果は実時間側の防壁のため、\texttt{result} は期限切れの握り潰しのため、
セッション型は手順の取り違えのためだった。
AI エージェントという題材の上に置くと、
それらが\textbf{同じひとつのプログラムの別々の壊れ方}を押さえていることが分かる。
言い換えれば、AI エージェントを型で書くために新しい理論は要らなかった ---
並行プログラムの型付けの問題が、そのまま現れているだけである。
\end{finding}

\subsection{実物の corpus はどうだったか}

これは作った例なので、実際に書かれていたプログラムでも確かめた。
処理系のリポジトリにある \texttt{aios\_user\_three\_role\_ai.aipl} は、
本節と同じ三役の構成を持つ --- ただし本稿の機構を入れる前に書かれたものである。

\begin{term}
[warn] line 43: now without a deadline waits forever
[warn] line 49: await without a deadline waits forever
[warn] line 52: await without a deadline waits forever
[warn] line 55: now without a deadline waits forever
OK
\end{term}

\textbf{モデルを待つ四つの箇所すべてに期限が無い}。
セッションの宣言はあり（このプログラムは \texttt{protocol\_define} を使っている）、
順序については静的検査を通っている。
つまり、書いた本人が意識していた部分は正しく、
意識していなかった部分 --- 「返ってこなかったらどうなるか」--- が
全部抜けていた。

\subsection{見えないところ}

正直に書いておくべき限界が二つある。

\paragraph{モデルの中身は型の外にある。}
\texttt{ai\_call\_with\_system} が返すのは \texttt{string} である。
「査読役は OK か NG を返す」という約束は、
プロンプトに書いてあるだけで型には無い。
効果と期限が押さえるのは\textbf{呼び出しの外側}であって、
モデルの出力の形ではない。
精緻化型（refinement type）を載せる余地はあるが、本稿では扱っていない。

\paragraph{三実装の mock が違う。}
本節のプログラムは\textbf{型検査だけ}を三実装で突き合わせている
（三つとも「問題なし」で一致）。
出力を比較していないのは、モデルを呼ばない代替（mock）の応答が
三実装で違うからである。
これは実装間の\textbf{未解消の差}であり、
モデルを呼ぶプログラムを出力で比較できるようにするには、
mock の応答そのものを仕様に含める必要がある。

% =====================================================================

% =====================================================================
\section{代表的なサンプルプログラム}
\label{sec:samples}

最新仕様の代表として 20 本を挙げる。
出力はすべて実行して得たものであり、
うち 18 本は\textbf{三つの処理系すべてで一致する}ことを確認している（\S\ref{sec:impl}）。
残る 2 本（\texttt{s11}・\texttt{s16}）は出力にノード名を含むため自動比較から外し、
内容を目視で照合した。
各サンプルに\textbf{示すもの・メリット・デメリット}を付す。

\subsection{s01 --- アクターの基本}
\begin{aipl}
class Greeter {
  var count = 0;
  method init(name) { print("hello, " ++ name); }   // new のとき自動で呼ばれる
  method tick() : unit !{mut, log} {
    count = count + 1;
    print("tick " ++ count);
  }
}
var g = new Greeter("AIPL");
send g.tick();
send g.tick();
\end{aipl}
\begin{term}
hello, AIPL
tick 1
tick 2
\end{term}
\textbf{示すもの}：クラス・フィールド・\texttt{init}・非同期送信・文字列連結。
\texttt{++} が連結で \texttt{+} は数値専用である。
\textbf{メリット}：状態を持つ並行部品が 10 行で書け、排他制御が要らない。
\textbf{デメリット}：\texttt{send} は結果を返さないので、
順序に依存する処理を書くと\textbf{送った順にしか保証がない}。

\subsection{s02 --- 三種の送信}
コードと出力は \S\ref{sec:abcl} に示した。
\textbf{メリット}：待つ・待たないを\textbf{呼ぶ側}が選べる。
同じメソッドを実時間側は \texttt{send} で、
エージェント側は \texttt{future} で使える。
\textbf{デメリット}：三種あるぶん、初学者はどれを使うか迷う。
\texttt{now} を安易に選ぶと待ちが連鎖する。

\subsection{s03 --- \texttt{reply} の規律}
\begin{aipl}
class Router {
  method route(k) : string {
    if (k > 0) { reply("positive"); } else { reply("non-positive"); }
  }
}
var r = new Router();
print(now r.route(5)  timeout 100 else "?");
print(now r.route(-1) timeout 100 else "?");
\end{aipl}
\begin{term}
positive
non-positive
\end{term}
\textbf{示すもの}：戻り値型を宣言すると、全経路での \texttt{reply} が義務になる。
\texttt{else} を消せば型エラーになる。
\textbf{メリット}：「返信を忘れて相手が待ち続ける」が静的に止まる。
\textbf{デメリット}：判定が構文の走査なので保守的である。
\texttt{while} の本体にしか \texttt{reply} が無い場合、0 回実行されうるとみなして拒否する。

\subsection{s04 --- 効果}
コードは \S\ref{sec:koba} に示した。
\begin{term}
1
1
n=1
\end{term}
\textbf{メリット}：\texttt{!\{\}} と書けたメソッドは純粋だと分かるので、
自由に複製・再配置できる。分散配置の判断が\textbf{型から機械的に決まる}。
\textbf{デメリット}：語彙が 8 個に固定で、利用者が増やせない。
また順序を持たないので「取得したら解放する」が書けない。

\subsection{s05 --- 期限}
コードと出力は \S\ref{sec:koba} に示した。
\textbf{メリット}：どんな相手でも待ちが有限で終わる。
実時間側の設計が「最悪でも $n$ ミリ秒」で組める。
\textbf{デメリット}：\texttt{else} の値が\textbf{正常値と区別できない}。
上の例の $-1$ は「失敗」を意味するが、型は \texttt{int} のままである。
失敗を型で表すには直和型が要る。

\subsection{s06 --- 実時間側の防壁}
\begin{aipl}
class Planner {                       // エージェント側
  method plan(g) : string !{ai, net} { reply(ai_call("plan " ++ g)); }
}
class Servo {                         // 実時間側。ai も net も持たない
  var p = new Planner();
  method step() : unit !{time, log} {
    send p.plan("grasp");             // 送るだけなら許される
    print("servo step");
  }
}
var s = new Servo();
send s.step();
\end{aipl}
\begin{term}
servo step
\end{term}
\textbf{示すもの}：本稿の看板。
\texttt{send} を \texttt{now p.plan(...)} に変えると、
\texttt{Servo.step} が \texttt{ai, net} を宣言していないので型エラーになる。
\textbf{メリット}：設計上の誤りが規約ではなく型検査で止まる。
\textbf{デメリット}：この保証は\textbf{注釈を書いた場合にだけ}効く。
注釈を省くと本体から推論されるので、誤りは誤りのまま通る。
防壁は「宣言する」という利用者の行為に依存している。

\subsection{s07 --- \texttt{select}}
\begin{aipl}
class Worker {
  method job(k) : string { reply("direct:" ++ k); }
  method serve() : unit !{log} {
    print("waiting");
    select {
      case job(k) -> { print("got " ++ k); reply("via select:" ++ k); }
      timeout 500 -> { print("timed out"); }
    }
  }
}
var w = new Worker();
send w.serve();
print(now w.job(1) timeout 800 else "?");
\end{aipl}
\begin{term}
waiting
got 1
via select:1
\end{term}
\textbf{示すもの}：\texttt{case} 本体の \texttt{reply} は、
囲む \texttt{serve} ではなく\textbf{受理した \texttt{job} への返信}になる。
\textbf{メリット}：待ち受けと処理が一箇所に書ける。
\textbf{デメリット}：\texttt{where} が無いので受理条件を状態で絞れない。
また \texttt{reply} の帰属が直感に反しやすく、被覆検査もここで別扱いになる。

\subsection{s08 --- 状態機械}
\begin{aipl}
class Door {
  var open = false;
  method state() : unit !{log} {
    if (open) { print("opened"); } else { print("closed"); }
  }
  method toggle() : unit !{mut, log} {
    if (open) { print("closing"); open = false; }
    else      { print("opening"); open = true; }
  }
}
var d = new Door();
send d.state();
send d.toggle();
send d.state();
\end{aipl}
\begin{term}
closed
opening
opened
\end{term}
\textbf{示すもの}：振る舞いの差し替えを、状態フィールドと分岐で書く。
アクター型が実行時に変わらないので、外から見た型が嘘にならない。
\textbf{メリット}：状態の数だけメソッドを増やさずに済む。
\textbf{デメリット}：状態が多いと \texttt{if} が深くなる。
状態遷移の構造そのものは、コードの見た目には現れない。

\subsection{s09 --- アクターを引数に渡す}
\begin{aipl}
class Device { method ping() : unit !{log} { print("pong"); } }
class Driver {
  method attach(d: Device) : unit !{log} {   // Device 以外を渡すと型エラー
    print("attached");
    send d.ping();
  }
}
var dev = new Device();
var drv = new Driver();
send drv.attach(dev);
\end{aipl}
\begin{term}
attached
pong
\end{term}
\textbf{示すもの}：引数に型注釈を書くとクラスが決まり、
\texttt{Device} 以外を渡すと \texttt{actor type mismatch} になる。
\textbf{メリット}：ドライバのような「相手を受け取って使う」部品が型安全に書ける。
\textbf{デメリット}：注釈は\textbf{必須}である。
省くと引数の型が決まらず、送信そのものが型エラーになる。
また注釈は\textbf{クラス名}であってインタフェースではないので、
「\texttt{ping} を持つ何か」を受け取ることはできない。

\subsection{s10 --- 三段のパイプライン}
\begin{aipl}
class Source { method get(i) : int !{} { reply(i * 10); } }
class Scale  {
  var src = new Source();
  method at(i) : int !{} { reply((now src.get(i) timeout 100 else 0) + 1); }
}
class Sink {
  var sc = new Scale();
  var total = 0;
  method run(n) : int !{mut, log} {
    var i = 0;
    while i < n do {
      total = total + (now sc.at(i) timeout 100 else 0);
      i = i + 1;
    }
    print("total=" ++ total);
    reply(total);
  }
}
var s = new Sink();
print("sum = " ++ (now s.run(4) timeout 1000 else -1));
\end{aipl}
\begin{term}
total=64
sum = 64
\end{term}
\textbf{示すもの}：フィールドに \texttt{new} でアクターを持ち、
\texttt{now} を二段重ねる実用的な形。
\textbf{メリット}：段ごとに独立したアクターなので、
後から段を差し替えたり分散配置したりできる。
\textbf{デメリット}：\texttt{now} の連鎖は\textbf{待ちの連鎖}である。
各段の期限が積み上がるので、外側の期限は内側の合計より大きくしなければならない。
上の例で外側を \texttt{timeout 100} にすると $-1$ になる。

\subsection{s11 --- メッシュへの配備}
コードと出力は \S\ref{sec:mesh} に示した。
\textbf{メリット}：どのノードに何を置けるかが\textbf{型から決まる}。
配置の判断が人手の規約でなくなる。
\textbf{デメリット}：ノードの方針は文字列で書く（\texttt{"mut, log, time"}）。
効果の語彙と綴りが一致しているかは実行時にしか分からない。

\subsection{s12 --- 失敗を型で表す}
コードと出力は \S\ref{sec:result} に示した。
\begin{term}
fast ok 42
\end{term}
\textbf{メリット}：期限切れを正常値として流せない。
\texttt{is\_ok} を書かずに使うと型エラーになる。
\textbf{デメリット}：待ちのたびに分岐が増える。
\texttt{else} を書く従来の形も残してあるので、
「必ず確かめる」を強制してはいない --- 選べるということは、選び損ねられるということでもある。

\subsection{s13 --- 資源の取得と解放}
コードと出力は \S\ref{sec:res} に示した。
\begin{term}
read = 42
\end{term}
\textbf{メリット}：デバイスを掴んだまま抜ける経路が静的に消える。
OS を書く目的に直結する。
\textbf{デメリット}：資源の名前が\textbf{文字列リテラル}である。
変数に入れた名前は追えない。
資源のあいだの順序は s19 で入れた。

\subsection{s14 --- 返信先の委譲}
コードと出力は \S\ref{sec:reply1st} に示した。
\begin{term}
ans = 42
\end{term}
\textbf{示すもの}：ABCL/1 から落とした機構の回復。
\textbf{メリット}：受付と処理を分けられる。
受けたアクターが自分で答えずに、専門のアクターへ丸ごと渡せる。
\textbf{デメリット}：線形性の検査は\textbf{メソッド内で閉じている}。
渡した先が答えるかどうかは、渡した先の検査に委ねられる。
各メソッドが局所的に規律を守ることで全体が成り立つ、という構成である。

\subsection{s15 --- 義務レベル}
コードと出力は \S\ref{sec:deadlock} に示した。
\begin{term}
level chain = 41
total = 41
\end{term}
\textbf{示すもの}：待ちが $1 \to 2 \to 3$ と一方向に流れる三段の連鎖。
\textbf{メリット}：全メソッドに注釈が付けば、循環待ちが\textbf{構成的に}起きない。
期限に頼らずに済む。
\textbf{デメリット}：レベルはメソッド単位なので、
同じクラスの別インスタンス同士の待ち（親が子を待つ木構造）も一律に弾かれる。
注釈の付け忘れは推論が埋めるので弱点ではなくなったが、
\textbf{注釈を書く意味}は残る --- コードが持つ構造より強い規律を先に宣言しておける。

\subsection{s16 --- ノードをまたぐ義務レベル}
コードと出力は \S\ref{sec:nodelevel} に示した。
\begin{term}
pos=3
drive = 3
\end{term}
\textbf{示すもの}：レベルが\textbf{メッシュの国境を越える}。
上位ノードは低いレベル、実時間ノードは高いレベル、という階層を宣言しておくと、
「実時間側が上位を待つ」という逆流が型で止まる。
\textbf{メリット}：分散デッドロックの主要な形（層をまたぐ待ちの循環）が消える。
\textbf{デメリット}：ノード単位なので粗い。
同じノードの中の待ちは、そのノードのレベル階層で別に見る必要がある。
またレベルを宣言していないノードは検査されない --- 効果の方針と同じ弱点である。

\subsection{s17 --- \texttt{select} の規律}
\begin{aipl}
class Worker3 {
  method job(k) : string { reply("done:" ++ k); }
  method serve() : unit !{log} {
    print("waiting");
    select {
      case job(k) -> { print("got " ++ k); reply("via select:" ++ k); }
      timeout 500 -> { print("timed out"); }
    }
  }
}
var w = new Worker3();
send w.serve();
print(now w.job(7) timeout 800 else "?");
\end{aipl}
\begin{term}
waiting
got 7
via select:7
\end{term}
\textbf{示すもの}：期限があり、待つメッセージ \texttt{job} を実際に誰かが送っている。
どちらか欠けると警告が出る。
\textbf{メリット}：\texttt{case} の綴り間違いと、書き忘れた送信が静的に出る。
\textbf{デメリット}：送信が動的だと見えない。
外部に公開しているプログラムでは検査そのものを行わない ---
\textbf{言語が世界の全部を見ていない}ことの現れである。

\subsection{s18 --- セッション型}
コードと出力は \S\ref{sec:session} に示した。
\begin{term}
pipeline = 41
\end{term}
\textbf{示すもの}：やり取りの順序を約束として書き、走らせる前に照合する。
\textbf{メリット}：手順の取り違えが実行前に出る。
実行時の検査と\textbf{同じ宣言}を使うので、二重に書く必要が無い。
\textbf{デメリット}：静的に見えるのは、
\texttt{protocol\_start} と \texttt{protocol\_end} が同じ並びにあるセッションだけ。
アクターをまたぐ場合は実行時の検査に任せる。

\subsection{s19 --- 資源への全体順序}
コードと出力は \S\ref{sec:resorder} に示した。
\begin{term}
w = 40
r = 21
\end{term}
\textbf{示すもの}：取得の入れ子から順序を読み取り、閉路が無いことを見る。
\textbf{メリット}：注釈が要らない。
逆順に取っている\textbf{二か所}が、そのままエラーに出る。
\textbf{デメリット}：名前がリテラルの \texttt{acquire} しか追えない。
そこだけは実行時の検査に任せる。

\subsection{s20 --- アクターをまたぐセッション}
コードと出力は \S\ref{sec:crosssession} に示した。
\begin{term}
cross = 41
\end{term}
\textbf{示すもの}：窓口役をひとつ呼び、残りの手順はその本体で進む ---
実際のプログラムの形。同期の呼び先を差し込んで一本の並びにする。
\textbf{メリット}：プログラムを変えずに、またいだセッションが検査できる。
\textbf{デメリット}：アクターは\textbf{引数で渡す}必要がある。
フィールドに持たせると、実行時の呼び名（\texttt{c\#f}）と
静的な呼び名（\texttt{f}）が食い違い、手順が役を指せない。

\subsection{20 本を通して見えること}

\begin{center}
\begin{tabular}{lll}
\toprule
& 効いている仕組み & 代価 \\
\midrule
s01--s03 & 型（\texttt{reply} の規律） & 判定が保守的 \\
s04, s06 & 効果 & 語彙が固定・注釈依存 \\
s05, s10 & 期限 & 待ちが連鎖する \\
s07--s09 & アクターの構造 & 受理条件・注釈・読みやすさ \\
s11 & 効果＋配備 & 方針が文字列 \\
s12 & 失敗を型で表す & 分岐が増える。従来の形も選べる \\
s13 & 順序つきの効果 & 名前が文字列。全体の順序が無い \\
s14 & 線形型 & 検査がメソッド内で閉じている \\
s15 & 義務レベル & 粒度がメソッド単位 \\
s16 & ノード階層 & 粒度がノード単位。宣言しなければ検査されない \\
s17 & select の規律 & 動的な送信・外部の送り手は見えない \\
s18 & セッション型 & 同じ並びで完結するセッションに限る \\
s19 & 資源への全体順序 & 名前がリテラルの acquire に限る \\
s20 & またぐセッション & 同期の呼び先しか差し込めない \\
\bottomrule
\end{tabular}
\end{center}

代価の欄を縦に読むと、埋め方の癖が見える ---
\textbf{s05 の「失敗が正常値と同じ型」は s12 が埋め、
s04 の「効果に順序が無い」は s13 が埋め、
s13 の「資源に順序が無い」は s19 が埋め、
s18 の「またげない」は s20 が埋めている}。
一つの機構が残した穴を次の機構が埋める、という連なりになっている。
そして埋めるたびに新しい代価が付く ---
本節の 20 本は、その代価を並べたものでもある。

三つの仕組みはそれぞれ独立に理解できるが、
\textbf{噛み合って初めて言えること}が s06 である。
効果だけでも期限だけでも、あの誤りは止まらない。

% =====================================================================

% =====================================================================
\section{三つの処理系と、揃っていることの確かめ方}
\label{sec:impl}

AIPL には実装が三つある --- 正となる OCaml 版（手書き 8,023 行）、
Python の CLI 実装（15,101 行）、ブラウザで動く JavaScript 実装（6,331 行）である。
三つ作る理由は二つある。
JS 版はブラウザで動くので処理系を配らずに試してもらえること、
そして\textbf{三つ揃えること自体が仕様の検査になる}ことである。

確かめ方は二本立てである。
\texttt{run\_all\_impls.sh} が\textbf{通るプログラムの出力が一致するか}を見て、
\texttt{run\_reject\_tests.sh} が\textbf{通ってはいけないプログラムが止まるか}を見る。
後者は\textbf{31 本}の「弾かれるべきプログラム」からなる。
前者だけでは足りない ---
検査が丸ごと抜け落ちていても、正しいプログラムの出力は一致するからである。

\subsection{現在の状態}

サンプル 20 本のうち\textbf{18 本が三実装で出力一致}する
（\texttt{s11} と \texttt{s16} は出力にノード名 \texttt{rt0/servo1} を含み、
比較スクリプトが区切りに使う \texttt{/} で割れてしまうため自動比較から外した。
内容は目視で一致を確認している）。
ガイドの 8 本も一致する。
弾かれるべき 31 本については、93 枠のうち\textbf{10 枠}が期待と食い違い、
その内訳は JS 版 7・Py 版 3 である
（\textbf{OCaml 版は 31 本すべてで期待どおりに振る舞う}）。
残りはいずれも「OCaml 版にはあるが他の実装に無い検査」であり、
言語仕様の穴ではなく追随の遅れである。

\subsection{仕様を変えるときの手順}

機構を一つ足すとき、次の順で進めている。

\begin{enumerate}
\item 最小の再現プログラムを書き、\textbf{実際の処理系に通す}。
\item 582 本の \texttt{.aipl} 全体に対する影響を測る ---
      \textbf{採用の前に}、何本が新たに落ちるかを数える。
\item その例を「弾かれるべきプログラム」として恒久的に登録する。
\item 三実装すべてに移植し、出力一致と拒否表を取り直す。
\end{enumerate}

第 2 段が要点である。
たとえば資源への全体順序は 582 本で新たに落ちるものが 0 本、
アクターをまたぐセッションの検査は旧版の型検査器と一件ずつ突き合わせて
判定の変わったファイルが 0 本であった。
\textbf{「厳しくしたが既存を壊していない」を数で言えるようにしておく}。

\begin{finding}
「三実装が揃っている」という主張は、
\textbf{正しいプログラムの出力について}しか成り立っていなかった。
「弾かれるべきプログラム」の側を作って初めて、
実装間の食い違いが見えるようになった ---
最初に表を取ったとき、11 枠が食い違っていた。
検査は、通す側からは見えない。
\end{finding}

\subsection{機械検証との距離}

正直に書いておくべきことがある。
機械検証（Coq/Rocq、\texttt{AIPLSoundness.v}、1121 行、\texttt{Print Assumptions} が公理なし）が
済んでいるのは、\textbf{効果も期限も持たない中核} $\mathrm{AIPL}^{-}$ である ---
保存定理・進行定理・型安全性・\texttt{method not understood} の不在・
非同期送信のデッドロック自由性が示されている。
効果と期限を加えた $\mathrm{AIPL}^{-2}$ 以降は紙の上の証明であり、
本稿の七つの規律は機械検証されていない。
\textbf{実装と機械検証の距離は、縮まるどころか広がっている}。

これは望ましくないが、避けがたくもある。
機構を足すたびに形式化を追いかけると、
「実装が何をしているか」を確かめる速度が落ちる。
現状は、\textbf{機械検証の代わりに「弾かれるべきプログラム」31 本}が
実装の規律を押さえている。
これは証明ではないが、\textbf{三つの実装すべてに同じ問いを投げられる}という利点がある ---
証明は一つのモデルについてしか言えない。

% =====================================================================
\section{考察}
\label{sec:discuss}

\subsection{採否の判断は一貫していたか}

表（\S\ref{sec:table}）を通して見ると、判断の軸は三つに整理できる。

\paragraph{軸 1 --- 反例が AIPL の言葉で出るか。}
交差型による高階モデル検査を落としたのは、検証力が足りないからではない。
反例が高階再帰スキームの言葉で返り、
「どのアクターがどのメッセージで詰まるか」にならないからである。
逆に義務レベルを採れたのは、推論したレベルがそのまま
「\texttt{Ctl2\#go (@7)} が \texttt{@1} を待っている」という形の証人になったからである。

\paragraph{軸 2 --- 証明の構造を壊さないか。}
\textit{express mode} と \texttt{where} ガードは、
どちらも\textbf{進行定理に枝を増やす}という一点で落としている。
逐次性が破れれば \texttt{mut} の意味が変わり、
ガードが増えれば「受理できるメッセージが無い」構成を扱う必要が出る。
機能の魅力ではなく、土台への影響で決めている。

\paragraph{軸 3 --- 実際に書かれているか。}
判断を数で決めた場面がいくつかある ---
「名前が動的な \texttt{acquire}」を課題から外したのは、
582 本を測って\textbf{一件も無かった}からである。
逆にアクターをまたぐセッションを追いかけたのは、
セッションを使う 13 本のうち\textbf{6 本}がまたいでいたからである。
\textbf{仕様の穴は、大きさではなく踏まれ方で優先順位が決まる}。

\subsection{型・効果・期限という三点セットの評価}

三つの仕組みの効き方には差がある。

\textbf{型}は最もよく効いている。
\texttt{reply} の規律とアクター型の区別は、書き間違いの多くをその場で止める。
代価は判定の保守性で、正しいのに拒否される場合がある。

\textbf{効果}は、効くときは非常に効くが\textbf{注釈依存}である。
書かなければ推論されるだけで、誤りは通る。
これは「効果を書かないと損をする」仕組みが無いためで、
たとえば公開するアクターには注釈を義務づけるといった運用が要る。
一方、メッシュ配備では \texttt{node\_allow} が\textbf{注釈を要求する側}に回るので、
国境をまたぐコードについては強制が効く。

\textbf{期限}は最も安上がりに強い性質を与える。
構文を一つ足すだけで進行定理の枝が消える。
ただしこれは失敗を消すのではなく\textbf{失敗の形を変える}だけである。
\texttt{result}$\langle\tau\rangle$ を入れて「変わった形の失敗」を型に出せるようにしたが、
\texttt{else} を書けば従来どおりなので、規律は書き手が選ぶことになる。

\subsection{足した機構に共通していたこと}

七つの規律のうち、後から足した四つ ---
義務レベル・\texttt{select} の規律・セッション型・資源への全体順序 ---
には、共通する性質があった。
\textbf{どれも新しい注釈を作っていない}。

\begin{itemize}
\item 義務レベルは、\textbf{待ちの辺}を読んだ（効果の伝播に使う辺がそのまま使えた）。
\item セッション型は、\textbf{実行時に既にあった宣言}を読んだ。
\item 資源への全体順序は、\textbf{取得の入れ子}を読んだ
      （$r$ を持ったまま $s$ を取ったなら $r < s$）。
\item \texttt{select} の規律は、\textbf{他の場所にある送信}を読んだ。
\end{itemize}

必要な情報は、たいていプログラムの中に既に書かれている。
型システムの仕事は、書かせることより\textbf{読み取ること}であった。

\begin{finding}
「アクターをまたぐから静的には追えない」という判断も、
実は\textbf{型の限界ではなく検査の範囲の狭さ}だった。
必要だったのは多者間セッション型でも委譲でもなく、
同期の呼び先の送信列を差し込むことと、名前の解決表である。
プログラムは一行も変えていない。

そして限界は、予想していたところ（またぐこと）ではなく\textbf{名前}に現れた ---
実行時はフィールドのアクターを持ち主で修飾した名前（\texttt{c\#f}）で呼ぶが、
静的な側はクラスの実体を一つに決められない。
\end{finding}

\subsection{「正」の実装という考え方の限界}

OCaml 版を「正」と決めているが、これは\textbf{どちらが正しいかを決める手続き}ではない。
三実装を突き合わせて食い違いが出たとき、
OCaml 版が正しいとは限らない ---
実際、突き合わせで見つかった実装のバグのうち、
半分以上が\textbf{正であるはずの OCaml 版}のものであった。
「正」は\textbf{仕様の置き場所}であって、正しさの根拠ではない。

根拠になるのは、\textbf{同じ問いを三つに投げられること}である。
機械検証は一つのモデルについてしか言えないが、
「弾かれるべきプログラム」は三つ全部に投げられる。

\subsection{残っている課題}

\begin{enumerate}
\item \textbf{非同期の呼び先をまたぐセッション}。
      \S\ref{sec:crosssession} の展開は同期の呼び先しか差し込めない。
      \texttt{send} / \texttt{future} の先で手順が進む場合は、
      順序が決まらないので実行時に任せている。
      多者間セッション型と委譲（線形な \texttt{session}$\langle S\rangle$ を
      引数で引き回す）に進めば追えるが、プログラムの書き換えが要る。
\item \textbf{フィールドに持ったアクターは役になれない}。
      実行時は持ち主で修飾した名前で呼ぶが、
      静的な側はクラスの実体を一つに決められない。
      手順が名前で役を指す設計そのものに由来する。
\item \textbf{資源がアクターの実体であるときの順序}。
      \S\ref{sec:resorder} の順序は\textbf{資源の名前}についてのものである。
      食事する哲学者のように資源が\textbf{同じクラスの別インスタンス}
      （\texttt{lo} / \texttt{hi}）である場合、
      義務レベルの粒度が \texttt{Class\#method} なので区別できない。
      実体ごとの順序には精緻化型 $\{i:\mathtt{int} \mid i < j\}$ か、
      フィールドによる順序宣言が要る。
\item \textbf{モデルの出力の形は型の外にある}。
      「査読役は OK か NG を返す」はプロンプトにしか書かれていない。
      効果と期限が押さえるのは\textbf{呼び出しの外側}である。
\item \textbf{モデルの mock が三実装で違う}。
      そのためモデルを呼ぶプログラムは型検査でしか突き合わせられない。
\item \textbf{メソッド内で起きた失敗の伝わり方}。
      呼び出し側の \texttt{now ... timeout ... else} が何を返すかが
      三実装で違う。\texttt{result}$\langle\tau\rangle$ に載せる道と揃える必要がある。
\item \textbf{$\mathrm{AIPL}^{-2}$ 以降の機械検証}。
      効果つき判断、期限、\texttt{future}$\langle\tau\,!\,\varepsilon\rangle$、
      そして七つの規律。
\item \textbf{実機のメッシュへの接続}。
      配備の輸送は現状ひとつのプロセス内で模している。
      Xinu ノード間の実際の転送に差し替える。
\end{enumerate}

% =====================================================================
\section{おわりに}

本稿は AIPL の現在の仕様と、二つの源流からの採否をまとめた。

ABCL/1 からは\textbf{三種の送信・内部の逐次性・\texttt{reply}・返信先の一級性}を採り、
\textit{express mode} と \texttt{where} ガードを落とした。
落とした二つは、どちらも\textbf{進行定理に枝を増やす}という一点で共通している。

小林らの型システムからは\textbf{資源使用解析（効果）・効果の順序・線形型・
デッドロック自由な型}を採り、サイズ型は\textbf{期限という粗い代用}に置き換えた。
交差型モデル検査・多相メソッド・完全な所有権型は落とした。
落とした三つは、\textbf{反例が言語の言葉で出ない}か、
\textbf{$\rho$ による \texttt{reply} の推論と両立しない}か、
\textbf{アクターの性質から既に成立している}かのいずれかである。

どちらの源流にも無いものとして、\textbf{セッション型・資源への全体順序・
メッシュ配備}を足した。

結果として、\textbf{待ちの宛先がクラスとして分かるか、
宛先ノードのレベルが宣言されている範囲では、デッドロック自由が型で保証される}。
その外側では、期限が失敗を有限時間に変換する。

\paragraph{分かったこと。}
AI エージェントを型で書くために、\textbf{新しい理論は要らなかった}。
1986 年の ABCL/1 と、1997 年以降の小林らの型システムと、
セッション型という三十年ぶんの蓄積が、そのまま当たる問題であった。
足りていなかったのは理論ではなく、
それらを\textbf{一つの言語に同時に載せること}である。

そして載せる作業を通して見えたのは、
\textbf{必要な情報はたいていプログラムの中に既に書かれている}ということであった ---
待ちの辺、取得の入れ子、実行時の宣言、他の場所にある送信。
後から足した規律は、どれも新しい注釈を作らずに済んだ。

\paragraph{仕様は、書くだけでは検査されない。}
走らせて、複数の実装で走らせて、そして\textbf{実際に使われているかを数えて}、
初めて検査される。

\subsection*{再現方法}

\begin{term}
# 本稿の20本を流す
cd ~/aios/abclcp
for f in docs/samples/paper/s*.aipl; do
  printf 'load %s\ncompile\n' "$PWD/$f" > /tmp/s.repl
  ./src/abclrepl_thread -q -f /tmp/s.repl
done

# 通るプログラムの出力一致 / 弾かれるべき31本の拒否（三実装）
sh tools/run_all_impls.sh
sh tools/run_reject_tests.sh
\end{term}

% =====================================================================
\begin{thebibliography}{99}

\bibitem{abcl1}
Akinori Yonezawa, Jean-Pierre Briot, Etsuya Shibayama.
\newblock Object-Oriented Concurrent Programming in ABCL/1.
\newblock \emph{OOPSLA 1986}, pp.~258--268.

\bibitem{abclbook}
Akinori Yonezawa (ed.).
\newblock \emph{ABCL: An Object-Oriented Concurrent System}. MIT Press, 1990.

\bibitem{abclf}
Kenjiro Taura, Satoshi Matsuoka, Akinori Yonezawa.
\newblock ABCL/f: A Future-Based Polymorphic Typed Concurrent
Object-Oriented Language. 1994.

\bibitem{abclr}
Takuo Watanabe, Akinori Yonezawa.
\newblock Reflection in an Object-Oriented Concurrent Language (ABCL/R).
\newblock \emph{OOPSLA 1988}.

\bibitem{hewitt}
Gul Agha.
\newblock \emph{Actors: A Model of Concurrent Computation in Distributed
Systems}. MIT Press, 1986.

\bibitem{kobayashi94}
Naoki Kobayashi, Akinori Yonezawa.
\newblock Type-Theoretic Foundations for Concurrent Object-Oriented Programming.
\newblock \emph{OOPSLA 1994}.

\bibitem{kptl96}
Naoki Kobayashi, Benjamin C. Pierce, David N. Turner.
\newblock Linearity and the Pi-Calculus. \emph{POPL 1996}.

\bibitem{kobadf}
Naoki Kobayashi.
\newblock A Partially Deadlock-Free Typed Process Calculus (1997) /
A New Type System for Deadlock-Free Processes (\emph{CONCUR 2006}).

\bibitem{kobagen01}
Naoki Kobayashi.
\newblock A Generic Type System for the Pi-Calculus. \emph{POPL 2001}.

\bibitem{kobares}
Naoki Kobayashi et al.
\newblock Resource Usage Analysis. 2001/2002/2005.

\bibitem{koba09}
Naoki Kobayashi, C.-H. Luke Ong.
\newblock A Type System Equivalent to the Modal Mu-Calculus Model Checking
of Higher-Order Recursion Schemes. 2009.

\bibitem{kobasize}
Naoki Kobayashi et al.
\newblock Sized Types and Parallel Complexity.

\bibitem{akka}
Lightbend.
\newblock Akka Typed --- \texttt{Behavior[T]} and behavior replacement.

\bibitem{otp}
Ericsson.
\newblock Erlang/OTP \texttt{gen\_statem} Behaviour.

\bibitem{selfpaper1}
児玉靖司.
\newblock AIPL: Actor based Intelligence Parallel Language --- 設計から実現.
\newblock 2026 年 8 月 21 日（第 1 版。形式的保証と穴の修正）.

\bibitem{selfpaper2}
児玉靖司.
\newblock 同上 第 2 版 --- ABCL/1 と小林の型システムから、何を採り何を落としたか.
\newblock 2026 年 8 月 22 日.

\bibitem{selfvsabcl1}
児玉靖司.
\newblock AIPL（ABCLc+）と ABCL/1 --- 機能とモデルの観点からの比較. 2026 年 7 月 30 日.

\bibitem{selffeatures}
児玉靖司.
\newblock AIPL の言語仕様の検討 --- ABCL/1 の機構の再導入と、
分散 AI エージェント／実時間 OS という目的から. 2026 年 7 月 30 日.

\bibitem{selfsound2}
児玉靖司.
\newblock AIPL の型健全性と型安全性（第 2 版）. 2026 年 7 月 30 日.

\bibitem{selfreply}
児玉靖司.
\newblock AIPL におけるメソッド戻り値型の推論 --- reply の引数から型は決まるか.
2026 年 7 月 29 日.

\bibitem{selfguide}
児玉靖司.
\newblock AIPL / ABCLc+ ユーザーズガイド. 2026 年 7 月 30 日.

\bibitem{selfkoba}
児玉靖司.
\newblock 小林研究の成果を AIPL に取り入れる場合の言語仕様 ---
可能性、メリット、デメリット. 2026 年 7 月 30 日.

\end{thebibliography}

\end{document}
